\chapter{Complexity — SAT}
%%% generated by the blueprint pipeline from TCSlib.Complexity.ClassNP.SAT
%%%%%%%%%%%%%%%%%%%%%%%%%%%%%%%%%%%%%%%%%%%%%%%%%%%%%%%%%%%%%%%%%%%%%%%%%%%%%%%
\section{Overview}

This module defines the languages $\mathrm{SAT}$ and $\mathrm{3SAT}$ over the audited
binary string encoding of CNF formulas, and proves the membership half of the
Cook--Levin plan: $\mathrm{SAT} \in \mathrm{NP}$ and $\mathrm{3SAT} \in \mathrm{NP}$
(the satisfying assignment is the certificate), together with the clause-splitting
reduction $\mathrm{SAT} \le_p \mathrm{3SAT}$ [AB09, \S2.3.1, Theorem~2.10,
Lemma~2.14]. Decoding of strings is total, with the empty (satisfiable, vacuously
3CNF) formula as fallback, so every non-well-formed string lies in both languages.
The bulk of the module is a private layer of concrete machine constructions: a
finite-state syntax scanner for the CNF grammar, a streaming formula evaluator,
a clause-splitting transform with fresh-variable allocation and its equisatisfiability
proof, and a modular streaming transducer (pairing, suffix extraction, append,
padding, call-stopping, and host-loop combinators) realizing the reduction in
polynomial time.

Throughout this chapter, a CNF formula over $\mathbb{N}$-indexed variables is a
list of clauses, each clause a list of literals $(v,b)$ with variable index
$v \in \mathbb{N}$ and polarity $b \in \{0,1\}$. The serialization of a formula
writes each clause as a leading $1$ followed by its serialized literals — a literal
$(v,b)$ becomes $1^{v+1}\,0\,b$ — and a closing $0$, with one final $0$ terminating
the formula. The pair encoding $\mathrm{pair}(u,v)$ of two words doubles every bit
of $u$, then appends the separator $01$ followed by $v$.

%%%%%%%%%%%%%%%%%%%%%%%%%%%%%%%%%%%%%%%%%%%%%%%%%%%%%%%%%%%%%%%%%%%%%%%%%%%%%%%
\section{Declarations}

\begin{lemma}[Linear time implies polynomial time]
\lean{Complexity.sat_pt_linear}
\label{Complexity.sat_pt_linear}
\leanok
\uses{Complexity.PolyTimeComputable, Complexity.polyTimeComputable_of_linear, Turing.FinTM, Turing.FinTM.ComputesFunInTime}
\difficulty{1}
Let $f$ be a function from binary strings to binary strings. If there exist a finite
multi-tape Turing machine over the binary alphabet and a constant $C$ such that the
machine computes $f$ within $C\cdot(n+1)$ steps on every input of length $n$, then
$f$ is polynomial-time computable.
\end{lemma}

\begin{lemma}[Constant word is polynomial-time computable]
\lean{Complexity.sat_pt_const}
\label{Complexity.sat_pt_const}
\leanok
\uses{Complexity.PolyTimeComputable, Complexity.polyTimeComputable_const}
\difficulty{1}
For every fixed binary word $w$, the constant function sending every binary string
to $w$ is polynomial-time computable.
\end{lemma}

\begin{lemma}[Conditional of polynomial-time functions]
\lean{Complexity.sat_pt_cond}
\label{Complexity.sat_pt_cond}
\leanok
\uses{Complexity.PolyTimeComputable, Complexity.polyTimeComputable_ite}
\difficulty{1}
Let $p$ be a Boolean predicate on binary strings and $f, g$ functions on binary
strings. If the single-bit function $x \mapsto [p(x)]$ and both $f$ and $g$ are
polynomial-time computable, then so is the function mapping $x$ to $f(x)$ when
$p(x)$ holds and to $g(x)$ otherwise.
\end{lemma}

\begin{lemma}[Conjunction of polynomial-time bits]
\lean{Complexity.sat_pt_and}
\label{Complexity.sat_pt_and}
\leanok
\uses{Complexity.PolyTimeComputable, Complexity.polyTimeComputable_and}
\difficulty{1}
Let $p$ and $q$ be Boolean predicates on binary strings. If the single-bit functions
$x \mapsto [p(x)]$ and $x \mapsto [q(x)]$ are both polynomial-time computable, then
so is $x \mapsto [p(x) \wedge q(x)]$.
\end{lemma}

\begin{lemma}[Composition with a machine correct on an image]
\lean{Complexity.sat_comp_on_image}
\label{Complexity.sat_comp_on_image}
\leanok
\uses{Turing.FinTM, Turing.FinTM.ComputesFunInTime, Turing.FinTM.ComputesInTime, Turing.FinTM.exists_comp_on_image}
\difficulty{1}
Let $M$ and $U$ be finite multi-tape Turing machines over the binary alphabet, let
$f, g$ be functions on binary strings and $T_1, T_2 : \mathbb{N} \to \mathbb{N}$ time
bounds. If $M$ computes $f$ within $T_1$, and for every input $x$ the machine $U$ run
on $f(x)$ halts within $T_2(|x|)$ steps with output $g(x)$, then some finite machine
computes $g$ within time $n \mapsto 2\,T_1(n) + T_2(n) + 2$.
\end{lemma}

\begin{definition}[The language SAT]
\lean{Complexity.SAT}
\label{Complexity.SAT}
\leanok
\uses{Std.Sat.CNF.Satisfiable, Std.Sat.CNF.decode}
The language $\mathrm{SAT}$ over the binary alphabet is the set of all strings
whose total CNF decoding is a satisfiable formula [AB09, Section~2.3.1]. Decoding
is total, with the empty (satisfiable) formula as fallback, so every
non-well-formed string belongs to $\mathrm{SAT}$.
\end{definition}

\begin{definition}[The language 3SAT]
\lean{Complexity.SAT3}
\label{Complexity.SAT3}
\leanok
\uses{Std.Sat.CNF.Satisfiable, Std.Sat.CNF.WidthAtMost, Std.Sat.CNF.decode}
The language $\mathrm{3SAT}$ over the binary alphabet is the set of all strings
whose decoded CNF formula has width at most three (every clause carries at most
three literal occurrences) and is satisfiable [AB09, Section~2.3.1]. The fallback
formula has no clauses, so non-well-formed strings belong to $\mathrm{3SAT}$ as
well.
\end{definition}

\begin{definition}[Assignment carried by a certificate]
\lean{Complexity.satAssignment}
\label{Complexity.satAssignment}
\leanok
Given a binary word $u$, the associated truth assignment
$\mathbb{N} \to \{0,1\}$ maps a variable index $v$ to the $v$-th bit of $u$, and to
$0$ (false) whenever $v \ge |u|$.
\end{definition}

\begin{lemma}[Certificate characterization of satisfiability]
\lean{Complexity.sat_certificate}
\label{Complexity.sat_certificate}
\leanok
\uses{Complexity.eval_congr_of_lt_numVars, Complexity.satAssignment, Std.Sat.CNF.Satisfiable, Std.Sat.CNF.decode, Std.Sat.CNF.numVars_decode_le}
\difficulty{4}
For every binary string $x$, the CNF formula decoded from $x$ is satisfiable if and
only if there exists a word $u$ of length exactly $|x| + 1$ such that the decoded
formula evaluates to true under the assignment that reads variable $v$ off the
$v$-th bit of $u$ (default false). [AB09, Theorem~2.10, membership]
\end{lemma}

\begin{lemma}[Split success equation]
\lean{Complexity.sat_split_some}
\label{Complexity.sat_split_some}
\leanok
\uses{Turing.solveSplit}
\difficulty{2}
If the certificate-split solver with parameters $(1,1)$, searching for a split of a
total length $N$, returns the value $i$, then $i + (i + 1) = N$.
\end{lemma}

\begin{lemma}[Split completeness]
\lean{Complexity.sat_split_exists}
\label{Complexity.sat_split_exists}
\leanok
\uses{Complexity.sat_split_some, Turing.solveSplit}
\difficulty{3}
If natural numbers $i$ and $N$ satisfy $i + (i + 1) = N$, then the certificate-split
solver with parameters $(1,1)$ on total length $N$ succeeds and returns exactly $i$;
in particular this covers $N = 1$, $i = 0$.
\end{lemma}

\begin{definition}[Boolean width-three test]
\lean{Complexity.satWidth}
\label{Complexity.satWidth}
\leanok
The Boolean test on a CNF formula that checks that every clause has at most three
literal occurrences; repeated literals count as distinct occurrences.
\end{definition}

\begin{lemma}[Width test agrees with the width predicate]
\lean{Complexity.satWidth_spec}
\label{Complexity.satWidth_spec}
\leanok
\uses{Complexity.satWidth, Std.Sat.CNF.WidthAtMost}
\difficulty{1}
For every CNF formula $\varphi$, the Boolean width scan returns true if and only if
$\varphi$ has width at most $3$, i.e.\ every clause of $\varphi$ has at most three
literals.
\end{lemma}

\begin{definition}[Total mathematical verifier]
\lean{Complexity.satVerdict}
\label{Complexity.satVerdict}
\leanok
\uses{Complexity.satAssignment, Complexity.satWidth, Std.Sat.CNF.decode, Turing.solveSplit}
Given a flag (plain SAT or 3SAT mode) and a binary string $z$, the verifier verdict:
split $|z|$ as $i + (i+1)$ and reject outright if no such split exists; otherwise
decode the first $i$ bits of $z$ to a formula $\varphi$ and read an assignment off
the remaining $i + 1$ bits, and accept exactly when (in 3SAT mode) $\varphi$ passes
the width-three test and (in both modes) $\varphi$ evaluates to true under that
assignment.
\end{definition}

\begin{definition}[Verifier languages]
\lean{Complexity.satVerifier}
\label{Complexity.satVerifier}
\leanok
\uses{Complexity.satVerdict}
For each mode flag, the language of all binary strings on which the total verifier
verdict accepts; these are the two verifier languages for the prescribed
certificate parameters $(1,1)$.
\end{definition}

\begin{lemma}[Verdict on a correctly sized concatenation]
\lean{Complexity.satVerdict_append}
\label{Complexity.satVerdict_append}
\leanok
\uses{Complexity.satAssignment, Complexity.satVerdict, Complexity.satWidth, Complexity.sat_split_exists, Std.Sat.CNF.decode}
\difficulty{2}
Let $x$ and $u$ be binary words with $|u| = |x| + 1$. Then the verifier verdict on
the concatenation $x \mathbin{+\!\!+} u$ equals the conjunction of the (mode-guarded)
width-three test of the formula decoded from $x$ with the evaluation of that formula
under the assignment read off $u$.
\end{lemma}

\begin{lemma}[SAT certificate equivalence]
\lean{Complexity.sat_verifier_equiv}
\label{Complexity.sat_verifier_equiv}
\leanok
\uses{Complexity.SAT, Complexity.satVerdict, Complexity.satVerdict_append, Complexity.satVerifier, Complexity.sat_certificate, Std.Sat.CNF.Satisfiable, Std.Sat.CNF.decode}
\difficulty{3}
A binary string $x$ belongs to $\mathrm{SAT}$ if and only if there exists a
certificate $u$ with $|u| = |x| + 1$ such that the concatenation
$x \mathbin{+\!\!+} u$ belongs to the plain-mode verifier language.
\end{lemma}

\begin{lemma}[3SAT certificate equivalence]
\lean{Complexity.sat3_verifier_equiv}
\label{Complexity.sat3_verifier_equiv}
\leanok
\uses{Complexity.SAT3, Complexity.satAssignment, Complexity.satVerdict, Complexity.satVerdict_append, Complexity.satVerifier, Complexity.satWidth, Complexity.satWidth_spec, Complexity.sat_certificate, Std.Sat.CNF.Satisfiable, Std.Sat.CNF.WidthAtMost, Std.Sat.CNF.decode}
\difficulty{4}
A binary string $x$ belongs to $\mathrm{3SAT}$ if and only if there exists a
certificate $u$ with $|u| = |x| + 1$ such that $x \mathbin{+\!\!+} u$ belongs to the
3SAT-mode verifier language. The width condition depends only on the instance $x$,
so the same certificate as for $\mathrm{SAT}$ suffices.
\end{lemma}

\begin{lemma}[Unary scanner partition]
\lean{Complexity.sat_takeTrues_repr}
\label{Complexity.sat_takeTrues_repr}
\leanok
\uses{Std.Sat.CNF.takeTrues}
\difficulty{3}
For every binary string $x$, if the leading-ones scanner returns the count $k$ and
suffix $s$, then $x = 1^{k} \mathbin{+\!\!+} s$: the scanner partitions its input
into the counted run of ones and the remaining suffix.
\end{lemma}

\begin{lemma}[Literal parsing reconstructs its input]
\lean{Complexity.sat_parseLit_repr}
\label{Complexity.sat_parseLit_repr}
\leanok
\uses{Complexity.sat_takeTrues_repr, Std.Sat.CNF.parseLit, Std.Sat.CNF.serializeLit}
\difficulty{3}
If the literal parser succeeds on a binary string $x$, returning a literal $\ell$
and a remainder $r$, then $x$ is exactly the serialization of $\ell$ followed by
$r$.
\end{lemma}

\begin{lemma}[Clause parsing reconstructs its input]
\lean{Complexity.sat_parseClause_repr}
\label{Complexity.sat_parseClause_repr}
\leanok
\uses{Complexity.sat_parseLit_repr, Std.Sat.CNF.parseClause, Std.Sat.CNF.parseLit, Std.Sat.CNF.serializeClause}
\difficulty{5}
If the fuel-indexed clause parser succeeds on a binary string $x$, returning a
clause $C$ and a remainder $r$, then $x$ is exactly the serialization of $C$
(including its terminator) followed by $r$.
\end{lemma}

\begin{lemma}[Formula parsing reconstructs its input]
\lean{Complexity.sat_parseClauses_repr}
\label{Complexity.sat_parseClauses_repr}
\leanok
\uses{Complexity.sat_parseClause_repr, Std.Sat.CNF.parseClause, Std.Sat.CNF.parseClauses, Std.Sat.CNF.serialize, Std.Sat.CNF.serializeClause}
\difficulty{5}
If the fuel-indexed formula parser succeeds on a binary string $x$, returning a
formula $\varphi$ and a remainder $r$, then $x$ is exactly the serialization of
$\varphi$ — every clause marker and the final terminator — followed by $r$.
\end{lemma}

\begin{lemma}[Exact parsing is inverse serialization]
\lean{Complexity.sat_parse_repr}
\label{Complexity.sat_parse_repr}
\leanok
\uses{Complexity.sat_parseClauses_repr, Std.Sat.CNF.parse, Std.Sat.CNF.serialize}
\difficulty{3}
If the exact-consumption CNF parser succeeds on a binary string $x$, returning the
formula $\varphi$, then $x$ equals the serialization of $\varphi$.
\end{lemma}

\begin{definition}[LL(1) grammar transition]
\lean{Complexity.satSyntaxStep}
\label{Complexity.satSyntaxStep}
\leanok
The transition function of the six-state CNF syntax automaton over $\{0,1\}$ is
defined on the positions formula, clause, unary index, polarity, exact end, and
error as follows. At formula level a $1$ leads to clause level and a $0$ to exact
end; at clause level a $1$ leads to the index position and a $0$ back to formula
level; the index position stays on $1$ and moves to polarity on $0$; polarity
returns to clause level on either bit; exact end and error move to error on every
bit.
\end{definition}

\begin{definition}[Residual grammar of the syntax scanner]
\lean{Complexity.satSyntaxSuffix}
\label{Complexity.satSyntaxSuffix}
\leanok
\uses{Std.Sat.CNF.serialize, Std.Sat.CNF.serializeClause}
For each of the six scanner positions, the residual grammar of that position is
the set of binary strings that may legally remain to be read: at formula level,
the serializations of CNF formulas; at clause level, a serialized clause followed
by a serialized formula; at index level, a unary run $1^v$ followed by $0$, a
polarity bit, a serialized clause, and a serialized formula; at polarity level,
one arbitrary bit followed by a serialized clause and a serialized formula; at
exact end, only the empty string; at the error position, no string at all.
\end{definition}

\begin{lemma}[One grammar bit per transition]
\lean{Complexity.satSyntaxSuffix_cons}
\label{Complexity.satSyntaxSuffix_cons}
\leanok
\uses{Complexity.satSyntaxStep, Complexity.satSyntaxSuffix, Std.Sat.CNF.serialize, Std.Sat.CNF.serializeClause, Std.Sat.CNF.serializeLit}
\difficulty{5}
For every scanner position $q$, bit $b$, and binary string $x$: the string
$b :: x$ lies in the residual grammar of $q$ if and only if $x$ lies in the residual
grammar of the position reached from $q$ on reading $b$.
\end{lemma}

\begin{lemma}[Acceptance at end of input]
\lean{Complexity.satSyntaxSuffix_nil}
\label{Complexity.satSyntaxSuffix_nil}
\leanok
\uses{Complexity.satSyntaxSuffix, Std.Sat.CNF.serialize, Std.Sat.CNF.serializeClause}
\difficulty{2}
The empty string lies in the residual grammar of a scanner position $q$ if and only
if $q$ is the exact-end position.
\end{lemma}

\begin{lemma}[The scan recognizes the residual grammar]
\lean{Complexity.satSyntaxSuffix_run}
\label{Complexity.satSyntaxSuffix_run}
\leanok
\uses{Complexity.satSyntaxStep, Complexity.satSyntaxSuffix, Complexity.satSyntaxSuffix_cons, Complexity.satSyntaxSuffix_nil}
\difficulty{3}
For every scanner position $q$ and binary string $x$: $x$ lies in the residual
grammar of $q$ if and only if folding the grammar transition over $x$ starting at
$q$ ends in the exact-end position.
\end{lemma}

\begin{definition}[Boolean syntax pass]
\lean{Complexity.satSyntax}
\label{Complexity.satSyntax}
\leanok
\uses{Complexity.satSyntaxStep}
The Boolean result of the complete syntax scan of a binary string: fold the
six-state grammar transition over the whole string starting at the formula position
and test whether the exact-end position is reached.
\end{definition}

\begin{lemma}[Syntax scan agrees with the parser]
\lean{Complexity.satSyntax_spec}
\label{Complexity.satSyntax_spec}
\leanok
\uses{Complexity.satSyntax, Complexity.satSyntaxSuffix_run, Complexity.sat_parse_repr, Std.Sat.CNF.parse, Std.Sat.CNF.parse_serialize, Std.Sat.CNF.serialize}
\difficulty{4}
For every binary string $x$, the Boolean syntax scan of $x$ returns true exactly
when the exact-consumption CNF parser succeeds on $x$. This holds on every string,
including the empty input, unfinished literals, and strings with trailing garbage.
\end{lemma}

\begin{definition}[Generic one-way scanner machine]
\lean{Complexity.satScanTM}
\label{Complexity.satScanTM}
\leanok
\uses{Turing.FinTM}
Given a finite state set $S$ with a transition $S \times \{0,1\} \to S$, a start
state, and an accepting predicate on $S$, the one-way scanner is the finite Turing
machine with no work tapes that reads its input left to right applying the
transition, and upon reaching the right boundary emits the single bit given by the
accepting predicate at the final state and halts.
\end{definition}

\begin{definition}[Scanner configuration]
\lean{Complexity.satScanCfg}
\label{Complexity.satScanCfg}
\leanok
\uses{Turing.Cfg}
For an input word $x$, a control state $q$, and a position $i \le |x|$, the
scanner configuration is the configuration of a zero-work-tape machine on input
$x$ in control state $q$, with the input head just before position $i$ and empty
output.
\end{definition}

\begin{lemma}[One scanner step]
\lean{Complexity.satScan_step}
\label{Complexity.satScan_step}
\leanok
\uses{Complexity.satScanCfg, Complexity.satScanTM, Turing.Action.apply, Turing.Cfg.ext_zero_tapes, Turing.MultiTapeTM.step, Turing.inputSymbolInner, Turing.moveInputPos_pos_of_ne_right}
\difficulty{4}
Let a one-way scanner read input $x$ in state $q$ at position $i < |x|$. One machine
step moves silently to the state obtained by applying the transition to the $i$-th
input bit, with the head advanced to position $i+1$.
\end{lemma}

\begin{lemma}[Complete suffix scan]
\lean{Complexity.satScan_run}
\label{Complexity.satScan_run}
\leanok
\uses{Complexity.satScanCfg, Complexity.satScanTM, Complexity.satScan_step, Turing.Action.apply, Turing.Cfg.inputSymbol, Turing.MultiTapeTM.runFrom, Turing.MultiTapeTM.runFrom_succ_eq_step, Turing.MultiTapeTM.step}
\difficulty{4}
Let the input decompose as $x = p \mathbin{+\!\!+} r$. Started in any state $q$ at
position $|p|$, the one-way scanner halts after exactly $|r| + 1$ steps with output
the single bit obtained by applying the accepting predicate to the fold of the
transition over $r$ from $q$; no earlier step emits.
\end{lemma}

\begin{lemma}[Exact time bound for scanners]
\lean{Complexity.satScan_computes}
\label{Complexity.satScan_computes}
\leanok
\uses{Complexity.satScanCfg, Complexity.satScanTM, Complexity.satScan_run, Turing.Cfg.ext_zero_tapes, Turing.FinTM.ComputesFunInTime, Turing.FinTM.computesInTime_iff, Turing.MultiTapeTM.initCfg}
\difficulty{3}
Let $S$ be a finite set with transition function
$\delta : S \times \{0,1\} \to S$, start state $s_0 \in S$, and accepting
predicate $\alpha$ on $S$. On every input $x$ of length $n$, the one-way scanner
built from these data halts within $n + 1$ steps, and its output is the
single bit $\alpha(\delta^*(s_0, x))$, where $\delta^*(s_0, x)$ is the state
reached by folding $\delta$ over $x$ from $s_0$.
\end{lemma}

\begin{lemma}[The syntax pass is a real machine]
\lean{Complexity.satSyntax_poly}
\label{Complexity.satSyntax_poly}
\leanok
\uses{Complexity.PolyTimeComputable, Complexity.satScanTM, Complexity.satScan_computes, Complexity.satSyntax, Complexity.satSyntaxStep}
\difficulty{2}
The single-bit function returning the Boolean result of the complete CNF syntax
scan is polynomial-time computable; the witnessing machine is a one-way finite-state
scanner, and the result concerns every input string, not merely serialized formulas.
\end{lemma}
\begin{definition}[Evaluation control states]
\lean{Complexity.SatEvalControl}
\label{Complexity.SatEvalControl}
\leanok
The control-state type of the streaming formula evaluator: either one of four
administrative states (capture return, input rewind, dispatch, certificate rewind),
or a semantic phase — one of five streaming positions (formula, clause, unary index,
polarity, rewind) — together with a formula truth bit, a clause truth bit, and a
flag indicating that the second copy of a doubled input bit is to be skipped.
\end{definition}

\begin{definition}[Semantic evaluation state]
\lean{Complexity.satEvalQ}
\label{Complexity.satEvalQ}
\leanok
\uses{Complexity.SatEvalControl}
The embedding of a streaming phase, two accumulator bits, and a skip flag as a
semantic control state of the evaluator, with the input poised at a doubled bit
unless the skip flag is set.
\end{definition}

\begin{definition}[Administrative evaluation action]
\lean{Complexity.satEvalAction}
\label{Complexity.satEvalAction}
\leanok
\uses{Complexity.SatEvalControl, Turing.Action, Turing.FinTM}
Let $M$ be a base machine with $k$ work tapes. Given an input move $m$, a
certificate-head move $d$, an optional output bit $b$, and an optional target
state $q$, the administrative evaluation action is the action of the evaluator
over $M$ that writes nothing, keeps the first $k$ work heads stationary, moves the
input head by $m$ and the last (certificate) head by $d$, emits $b$ when present,
and switches to $q$ (halting when $q$ is absent).
\end{definition}

\begin{definition}[The streaming formula evaluator]
\lean{Complexity.satEvalTM}
\label{Complexity.satEvalTM}
\leanok
\uses{Complexity.SatEvalControl, Complexity.satEvalAction, Complexity.satEvalQ, Turing.FinTM, Turing.FinTM.controlAction, Turing.captureAction}
Given a certificate-extractor machine $M$ with $k$ work tapes, the evaluator is a
finite machine with $k + 1$ work tapes that first simulates $M$, capturing its
output (the certificate) on the fresh last tape, then rewinds the native input and
the certificate buffer, and finally streams a previously validated doubled
serialized formula, evaluating it clause by clause against the certificate: unary
index walks advance the certificate head, the polarity bit is compared with the
certificate cell, and clause and formula truth bits are accumulated in finite
control. No physical output occurs before the single final verdict bit.
[AB09, Theorem~2.10, membership]
\end{definition}

\begin{definition}[Canonical evaluation configuration]
\lean{Complexity.satEvalCfg}
\label{Complexity.satEvalCfg}
\leanok
\uses{Complexity.SatEvalControl, Complexity.satEvalTM, Turing.Cfg, Turing.FinTM, Turing.FinTM.bufferTape}
The canonical configuration of the evaluator during its streaming phase: the first
$k$ work tapes hold a saved configuration of the extractor machine, the last tape
holds the immutable certificate word $u$ with its head at a stated integer offset,
the control is a stated semantic state, the input head sits before a stated
position, and the output is empty.
\end{definition}

\begin{lemma}[Input symbol of the canonical configuration]
\lean{Complexity.satEvalCfg_input}
\label{Complexity.satEvalCfg_input}
\leanok
\uses{Complexity.SatEvalControl, Complexity.satEvalCfg, Turing.Cfg, Turing.Cfg.inputSymbol, Turing.FinTM, Turing.FinTM.inputSymbol_at}
\difficulty{1}
In the canonical evaluation configuration at input position $i$ on input $w$, the
symbol read by the input head is the optional $i$-th entry of $w$; it does not
depend on the saved extractor bank.
\end{lemma}

\begin{lemma}[Certificate cell under the last head]
\lean{Complexity.satEvalCfg_work}
\label{Complexity.satEvalCfg_work}
\leanok
\uses{Complexity.SatEvalControl, Complexity.satEvalCfg, Turing.Cfg, Turing.Cfg.workTapeSymbols, Turing.FinTM, Turing.FinTM.bufferTape}
\difficulty{1}
Let $u$ be a certificate word and $j \in \mathbb{Z}$. In every canonical
evaluation configuration with certificate $u$ and certificate head at offset $j$,
the symbol read by the last work head is the cell of the buffer tape of $u$ at
offset $j$, that is, the $j$-th bit of $u$ when $0 \le j < |u|$ and the blank
otherwise.
\end{lemma}

\begin{lemma}[Administrative action on canonical configurations]
\lean{Complexity.satEvalAction_apply}
\label{Complexity.satEvalAction_apply}
\leanok
\uses{Complexity.SatEvalControl, Complexity.satEvalAction, Complexity.satEvalCfg, Turing.Action.apply, Turing.Cfg, Turing.FinTM, Turing.moveInputPos}
\difficulty{3}
Applying a non-emitting administrative evaluator action to a canonical evaluation
configuration yields the canonical configuration with the prescribed new control
state, the input position moved as stated, and the certificate head displaced by
the stated amount; the saved extractor bank and the empty output are preserved.
\end{lemma}

\begin{lemma}[Silent skip of a doubled bit]
\lean{Complexity.satEval_skip}
\label{Complexity.satEval_skip}
\leanok
\uses{Complexity.satEvalAction, Complexity.satEvalAction_apply, Complexity.satEvalCfg, Complexity.satEvalQ, Complexity.satEvalTM, Turing.Action.apply, Turing.Cfg, Turing.FinTM, Turing.MultiTapeTM.step, Turing.moveInputPos_pos_of_ne_right}
\difficulty{3}
Let $w$ be an input and $i < |w|$. From a canonical evaluation configuration at
input position $i$ in a semantic state with phase $q$, accumulator bits $a$ and
$c$, and the skip flag set, one evaluator step yields the canonical configuration
at position $i + 1$ in the same phase with the same accumulator bits and the skip
flag cleared; the certificate head, the saved bank, and the empty output are
unchanged.
\end{lemma}

\begin{lemma}[Consuming a doubled bit]
\lean{Complexity.satEval_double}
\label{Complexity.satEval_double}
\leanok
\uses{Complexity.satEvalAction, Complexity.satEvalAction_apply, Complexity.satEvalCfg, Complexity.satEvalCfg_input, Complexity.satEvalQ, Complexity.satEvalTM, Complexity.satEval_skip, Turing.Action.apply, Turing.Cfg, Turing.Cfg.workTapeSymbols, Turing.FinTM, Turing.MultiTapeTM.runFrom, Turing.MultiTapeTM.step, Turing.moveInputPos_pos_of_ne_right}
\difficulty{4}
Suppose the evaluator's transition at a semantic state on the current input bit is
the administrative action moving right, displacing the certificate head by $d$, and
entering a new semantic state with the skip flag set. Then two machine steps — the
semantic transition followed by its silent skip — consume the doubled input bit:
the input position advances by two, the certificate head by $d$, and the control
reaches the new semantic state with the skip flag cleared. The certificate head
moves only on the first of the two physical transitions.
\end{lemma}

\begin{lemma}[Certificate rewind]
\lean{Complexity.satEval_rewind}
\label{Complexity.satEval_rewind}
\leanok
\uses{Complexity.SatEvalControl, Complexity.satEvalAction, Complexity.satEvalAction_apply, Complexity.satEvalCfg, Complexity.satEvalCfg_work, Complexity.satEvalTM, Turing.Action.apply, Turing.Cfg, Turing.FinTM, Turing.FinTM.bufferTape, Turing.FinTM.bufferTape_left, Turing.MultiTapeTM.runFrom, Turing.MultiTapeTM.runFrom_succ_eq_step, Turing.MultiTapeTM.runFrom_zero, Turing.MultiTapeTM.step, Turing.moveInputPos_zero}
\difficulty{5}
Let $q$ be an evaluator control state whose transition moves the certificate head
left while its cell is occupied and otherwise steps right into a destination state.
For every $n \le |u|$, starting from the canonical configuration with certificate
head at offset $n - 1$, exactly $n + 1$ steps return the head to offset $0$ in the
destination state, preserving all native input and output fields.
\end{lemma}

\begin{definition}[Bit doubling]
\lean{Complexity.satBits}
\label{Complexity.satBits}
\leanok
The word operation that doubles every bit of a binary string: each bit $b$ is
replaced by the two bits $b\,b$. This is the prefix format used by the native pair
encoding.
\end{definition}

\begin{lemma}[Length of a doubled word]
\lean{Complexity.satBits_length}
\label{Complexity.satBits_length}
\leanok
\uses{Complexity.satBits}
\difficulty{2}
For every binary string $x$, the doubled word has length exactly $2\,|x|$.
\end{lemma}

\begin{lemma}[Unary index walk]
\lean{Complexity.satEval_index}
\label{Complexity.satEval_index}
\leanok
\uses{Complexity.satBits, Complexity.satBits_length, Complexity.satEvalCfg, Complexity.satEvalQ, Complexity.satEvalTM, Complexity.satEval_double, Turing.Cfg, Turing.FinTM, Turing.MultiTapeTM.runFrom, Turing.MultiTapeTM.runFrom_add, Turing.MultiTapeTM.runFrom_zero}
\difficulty{5}
Suppose the evaluator's input contains, at the current position, the doubling of a
unary run $1^n$. In the index phase, exactly $2n$ steps consume the doubled run,
advancing the input position by $2n$ and the certificate head by $n$, leaving the
phase and both accumulator bits unchanged. (The first one of a literal is consumed
by the clause phase, so this scan counts the variable index itself.)
\end{lemma}

\begin{lemma}[Literal walk and rewind]
\lean{Complexity.satEval_literal}
\label{Complexity.satEval_literal}
\leanok
\uses{Complexity.satAssignment, Complexity.satBits, Complexity.satBits_length, Complexity.satEvalCfg, Complexity.satEvalCfg_work, Complexity.satEvalQ, Complexity.satEvalTM, Complexity.satEval_double, Complexity.satEval_index, Complexity.satEval_rewind, Std.Sat.CNF.serializeLit, Turing.Cfg, Turing.FinTM, Turing.FinTM.bufferTape, Turing.FinTM.bufferTape_nat, Turing.MultiTapeTM.runFrom, Turing.MultiTapeTM.runFrom_add}
\difficulty{6}
Suppose the evaluator's input contains, at the current position, the doubling of a
serialized literal $(v, b)$ with $v < |u|$ for the certificate $u$. Starting in the
clause phase with clause bit $c$, exactly $3v + 8$ steps process the literal: the
input advances by $2(v + 3)$, the certificate head walks to position $v$, the
polarity is compared with the $v$-th certificate bit, the clause bit becomes
$c \vee (u_v = b)$, and the certificate head returns to $0$.
\end{lemma}

\begin{lemma}[Clause evaluation pass]
\lean{Complexity.satEval_clause}
\label{Complexity.satEval_clause}
\leanok
\uses{Complexity.satAssignment, Complexity.satBits, Complexity.satBits_length, Complexity.satEvalCfg, Complexity.satEvalQ, Complexity.satEvalTM, Complexity.satEval_double, Complexity.satEval_literal, Std.Sat.CNF.serializeClause, Std.Sat.CNF.serializeLit, Turing.Cfg, Turing.FinTM, Turing.MultiTapeTM.runFrom, Turing.MultiTapeTM.runFrom_add}
\difficulty{5}
Let $u$ be the certificate word on the evaluator's last tape and $C$ a clause all
of whose variable indices are below $|u|$, with serialization $\mathrm{ser}(C)$.
Suppose the evaluator's input is a word $p$, followed by the doubling of
$\mathrm{ser}(C)$, followed by a word $r$. Started at input position $|p|$ in the
clause phase with formula bit $a$, clause bit $c$, and certificate head at $0$,
there is a number $t \le 3\,|\mathrm{ser}(C)|$ of steps after which the evaluator
is in the formula phase with formula bit $a \wedge (c \vee C(u))$, where $C(u)$ is
the truth value of $C$ under the assignment read off $u$, with the clause bit
reset, input position $|p| + 2\,|\mathrm{ser}(C)|$, and certificate head at $0$.
No output bit is emitted by this pass, even for an empty or false clause.
\end{lemma}

\begin{lemma}[Formula evaluation pass]
\lean{Complexity.satEval_formula}
\label{Complexity.satEval_formula}
\leanok
\uses{Complexity.satAssignment, Complexity.satBits, Complexity.satBits_length, Complexity.satEvalAction, Complexity.satEvalCfg, Complexity.satEvalCfg_input, Complexity.satEvalQ, Complexity.satEvalTM, Complexity.satEval_clause, Complexity.satEval_double, Std.Sat.CNF.serialize, Std.Sat.CNF.serializeClause, Turing.Cfg, Turing.Cfg.inputSymbol, Turing.FinTM, Turing.MultiTapeTM.runFrom, Turing.MultiTapeTM.runFrom_add, Turing.MultiTapeTM.runFrom_succ_eq_step, Turing.MultiTapeTM.runFrom_zero, Turing.MultiTapeTM.step}
\difficulty{5}
Let $u$ be the certificate word on the evaluator's last tape and $\varphi$ a CNF
formula all of whose variable indices are below $|u|$, with serialization
$\mathrm{ser}(\varphi)$. Suppose the evaluator's input is a word $p$, followed by
the doubling of $\mathrm{ser}(\varphi)$, followed by a word $r$. Started at input
position $|p|$ in the formula phase with formula bit $a$, clause bit false, and
certificate head at $0$, there is a number $t \le 3\,|\mathrm{ser}(\varphi)|$ of
steps after which the evaluator is halted and its output is the single bit
$a \wedge \varphi(u)$, where $\varphi(u)$ is the truth value of $\varphi$ under the
assignment read off $u$; only the final formula terminator emits.
\end{lemma}

\begin{lemma}[First halting time]
\lean{Complexity.sat_first_halt}
\label{Complexity.sat_first_halt}
\leanok
\uses{Turing.FinTM, Turing.FinTM.ComputesInTime, Turing.FinTM.computesInTime_iff, Turing.MultiTapeTM.initCfg, Turing.MultiTapeTM.runFrom, Turing.MultiTapeTM.runFrom_add, Turing.MultiTapeTM.runFrom_of_halt}
\difficulty{4}
Let a finite multi-tape machine $M$ compute output $y$ on input $x$ within $T$
steps. Then there is a time $t \le T$ such that $M$ is still live at every time
before $t$, halts exactly at time $t$, and its output at time $t$ is $y$: the first
halting transition can be selected while retaining the final emission.
\end{lemma}

\begin{lemma}[Evaluator startup seam]
\lean{Complexity.satEval_start}
\label{Complexity.satEval_start}
\leanok
\uses{Complexity.satEvalAction, Complexity.satEvalAction_apply, Complexity.satEvalCfg, Complexity.satEvalQ, Complexity.satEvalTM, Complexity.satEval_rewind, Complexity.sat_first_halt, Turing.Action.apply, Turing.Cfg, Turing.Cfg.init, Turing.FinTM, Turing.FinTM.ComputesInTime, Turing.FinTM.bufferTape, Turing.FinTM.timed_rewind, Turing.MultiTapeTM.initCfg, Turing.MultiTapeTM.runFrom, Turing.MultiTapeTM.runFrom_add, Turing.MultiTapeTM.runFrom_succ_eq_step, Turing.MultiTapeTM.step, Turing.captureCfg, Turing.capture_run, Turing.moveInputPos_zero}
\difficulty{6}
Let the extractor machine $M$ compute the certificate $u$ on input $w$ within $T$
steps. Then within $T + |w| + |u| + 5$ steps from its initial configuration, the
evaluator built over $M$ reaches the canonical evaluation configuration: some saved
bank of $M$, the certificate $u$ on the last tape with head at $0$, the input head
at $0$, formula phase with formula bit true, and empty output.
\end{lemma}

\begin{lemma}[One uniform linear-time evaluator]
\lean{Complexity.satEval_computes}
\label{Complexity.satEval_computes}
\leanok
\uses{Complexity.satAssignment, Complexity.satBits, Complexity.satBits_length, Complexity.satEvalTM, Complexity.satEval_formula, Complexity.satEval_start, Std.Sat.CNF.serialize, Turing.FinTM, Turing.FinTM.ComputesInTime, Turing.FinTM.ComputesInTime.mono, Turing.FinTM.computesFunInTime_pairSnd, Turing.FinTM.computesInTime_iff, Turing.MultiTapeTM.runFrom_add, Turing.pairDecode_pairEncode, Turing.pairEncode}
\difficulty{5}
There exist a single finite machine $E$ and a constant $A$ such that for every CNF
formula $\varphi$ and every certificate word $u$ covering all variables of
$\varphi$ (each literal index below $|u|$), the machine $E$, run on the pair
encoding of the serialization of $\varphi$ with $u$, halts within
$A \cdot (\ell + 1)$ steps — $\ell$ the length of that paired input — and outputs
the single bit giving the truth value of $\varphi$ under the assignment read off
$u$.
\end{lemma}
\begin{definition}[Catalog split of a verifier input]
\lean{Complexity.satSplit}
\label{Complexity.satSplit}
\leanok
\uses{Turing.pairEncode, Turing.solveSplit}
Given a binary string $z$: if $|z|$ splits as $i + (i + 1)$ for some $i$, return
the pair encoding of the first $i$ bits with the remaining $i + 1$ bits; otherwise
return the empty word as failure result.
\end{definition}

\begin{definition}[Recovered instance]
\lean{Complexity.satInstance}
\label{Complexity.satInstance}
\leanok
\uses{Complexity.satSplit, Turing.pairDecode}
For a binary string $z$, the recovered instance of $z$ is the first component of
the decoded catalog split of $z$, and is the empty word when the split or its
decoding fails.
\end{definition}

\begin{definition}[Recovered certificate]
\lean{Complexity.satWitness}
\label{Complexity.satWitness}
\leanok
\uses{Complexity.satSplit, Turing.pairDecode}
For a binary string $z$, the recovered certificate of $z$ is the second component
of the decoded catalog split of $z$ (the certificate region), and is the empty
word when the split or its decoding fails.
\end{definition}

\begin{definition}[Split validity]
\lean{Complexity.satSplitValid}
\label{Complexity.satSplitValid}
\leanok
\uses{Complexity.satSplit, Turing.pairDecode}
The Boolean test that the catalog split of a string succeeds and decodes as a pair;
in particular this is false for every string of even length.
\end{definition}

\begin{definition}[Full validity guard]
\lean{Complexity.satGood}
\label{Complexity.satGood}
\leanok
\uses{Complexity.satInstance, Complexity.satSplitValid, Complexity.satSyntax}
For a binary string $z$, the full validity guard of $z$ is the Boolean
conjunction of the split validity of $z$ with the complete CNF syntax scan of the
instance recovered from $z$; in particular it is false whenever the split of $z$
fails.
\end{definition}

\begin{definition}[Safe evaluator request]
\lean{Complexity.satSafe}
\label{Complexity.satSafe}
\leanok
\uses{Complexity.satGood, Complexity.satSplit, Std.Sat.CNF.serialize, Turing.pairEncode}
The safe input handed to the evaluator: the catalog split of the string when the
full validity guard holds, and otherwise the pair encoding of the serialized empty
formula with the empty certificate. Thus malformed formulas are replaced before the
evaluator runs, and no semantic rejection can precede validation.
\end{definition}

\begin{definition}[Safe evaluation value]
\lean{Complexity.satSafeValue}
\label{Complexity.satSafeValue}
\leanok
\uses{Complexity.satAssignment, Complexity.satGood, Complexity.satInstance, Complexity.satWitness, Std.Sat.CNF.decode}
The Boolean value of the evaluation after the syntax pass: when the full validity
guard holds, the truth value of the decoded instance under the assignment read off
the recovered certificate; otherwise true (the fallback empty formula evaluates to
true).
\end{definition}

\begin{lemma}[Literals lie below the variable bound]
\lean{Complexity.sat_literal_lt_numVars}
\label{Complexity.sat_literal_lt_numVars}
\leanok
\uses{Std.Sat.CNF.numVars}
\difficulty{2}
Let $\varphi$ be a CNF formula, $C$ a clause of $\varphi$, and $\ell$ a literal of
$C$. Then the variable index of $\ell$ is strictly below the variable bound of
$\varphi$ (the maximum over all occurring indices $v$ of $v + 1$).
\end{lemma}

\begin{lemma}[Safe requests are covered serializations]
\lean{Complexity.satSafe_spec}
\label{Complexity.satSafe_spec}
\leanok
\uses{Complexity.satAssignment, Complexity.satGood, Complexity.satInstance, Complexity.satSafe, Complexity.satSafeValue, Complexity.satSplit, Complexity.satSplitValid, Complexity.satSyntax_spec, Complexity.satWitness, Complexity.sat_literal_lt_numVars, Complexity.sat_parse_repr, Complexity.sat_split_some, Std.Sat.CNF.decode, Std.Sat.CNF.numVars, Std.Sat.CNF.numVars_decode_le, Std.Sat.CNF.parse, Std.Sat.CNF.serialize, Turing.pairDecode, Turing.pairDecode_pairEncode, Turing.pairEncode, Turing.solveSplit}
\difficulty{5}
For every binary string $z$ there exist a CNF formula $\varphi$ and a word $u$ such
that the safe evaluator request on $z$ is exactly the pair encoding of the
serialization of $\varphi$ with $u$, every variable index occurring in $\varphi$ is
below $|u|$, and the truth value of $\varphi$ under the assignment read off $u$
equals the safe evaluation value of $z$. This holds also on failed splits and
failed parses.
\end{lemma}

\begin{lemma}[The validation pipeline is polynomial-time]
\lean{Complexity.sat_pipeline_poly}
\label{Complexity.sat_pipeline_poly}
\leanok
\uses{Complexity.PolyTimeComputable, Complexity.PolyTimeComputable.comp, Complexity.polyTimeComputable_and, Complexity.polyTimeComputable_const, Complexity.polyTimeComputable_ite, Complexity.polyTimeComputable_of_linear, Complexity.satInstance, Complexity.satSafe, Complexity.satSplit, Complexity.satSplitValid, Complexity.satSyntax, Complexity.satSyntax_poly, Turing.FinTM.computesFunInTime_pairFst, Turing.FinTM.computesFunInTime_pairValid, Turing.FinTM.computesFunInTime_splitSolve}
\difficulty{3}
The following four functions on binary strings are polynomial-time computable: the
single-bit split-validity test, the instance projection, the single-bit syntax scan
of the projected instance, and the safe evaluator request assembly.
\end{lemma}

\begin{lemma}[The safe evaluation value is polynomial-time]
\lean{Complexity.satSafeValue_poly}
\label{Complexity.satSafeValue_poly}
\leanok
\uses{Complexity.PolyTimeComputable, Complexity.satEval_computes, Complexity.satSafe, Complexity.satSafeValue, Complexity.satSafe_spec, Complexity.sat_pipeline_poly, Turing.FinTM.ComputesInTime, Turing.FinTM.ComputesInTime.mono, Turing.FinTM.computesInTime_iff, Turing.FinTM.exists_comp_on_image, Turing.MultiTapeTM.output_length_le}
\difficulty{5}
The single-bit function sending a binary string $z$ to its safe evaluation value is
polynomial-time computable: the uniform evaluator is applied only on safe requests,
which are certified serializations with covering certificates.
\end{lemma}

\begin{lemma}[The SAT verdict is polynomial-time]
\lean{Complexity.satVerdict_false_poly}
\label{Complexity.satVerdict_false_poly}
\leanok
\uses{Complexity.PolyTimeComputable, Complexity.polyTimeComputable_const, Complexity.polyTimeComputable_ite, Complexity.satGood, Complexity.satInstance, Complexity.satSafeValue, Complexity.satSafeValue_poly, Complexity.satSplit, Complexity.satSplitValid, Complexity.satSyntax_spec, Complexity.satVerdict, Complexity.satWitness, Complexity.sat_pipeline_poly, Std.Sat.CNF.decode, Std.Sat.CNF.fallback, Std.Sat.CNF.parse, Turing.pairDecode, Turing.pairDecode_pairEncode, Turing.solveSplit}
\difficulty{4}
The single-bit function computing the plain-mode (SAT) verifier verdict on a binary
string is polynomial-time computable: failed splits are rejected, failed parses
take the accepting fallback branch, and all remaining inputs go through the safe
evaluation pipeline.
\end{lemma}

\begin{lemma}[A computed verdict decides its language]
\lean{Complexity.satVerifier_of_poly}
\label{Complexity.satVerifier_of_poly}
\leanok
\uses{Complexity.P, Complexity.PolyTimeComputable, Complexity.mem_P_iff, Complexity.satVerdict, Complexity.satVerifier, Turing.MultiTapeTM.indicator}
\difficulty{3}
For either mode flag: if the single-bit function computing the verifier verdict is
polynomial-time computable, then the corresponding verifier language belongs to the
class $\mathrm{P}$.
\end{lemma}

\begin{theorem}[SAT is in NP]
\lean{Complexity.SAT_mem_NP}
\label{Complexity.SAT_mem_NP}
\leanok
\uses{Complexity.NP, Complexity.SAT, Complexity.satVerdict_false_poly, Complexity.satVerifier, Complexity.satVerifier_of_poly, Complexity.sat_verifier_equiv}
\difficulty{2}
The language $\mathrm{SAT}$ — binary strings whose total CNF decoding is
satisfiable — belongs to $\mathrm{NP}$: a satisfying assignment, encoded as a
certificate of exactly $|x| + 1$ bits, is verified in polynomial time.
[AB09, Theorem~2.10, membership]
\end{theorem}

\begin{definition}[Saturated width counter]
\lean{Complexity.satWidthCap}
\label{Complexity.satWidthCap}
\leanok
For every $n \in \mathbb{N}$, the saturated width counter of $n$ is $\min(n, 3)$,
viewed as one of the four values $0, 1, 2, 3$; the width scanner keeps this
counter in finite control.
\end{definition}

\begin{definition}[Width-pass transition]
\lean{Complexity.satWidthStep}
\label{Complexity.satWidthStep}
\leanok
\uses{Complexity.satSyntaxStep, Complexity.satWidthCap}
The transition of the separate width pass, which is run only after the complete
syntax pass, acts on triples consisting of a position of the six-state syntax
automaton, a saturated literal counter in $\{0,1,2,3\}$, and a goodness flag. At
formula level it follows the syntax automaton and resets the counter; at clause
level a $1$ (the start of a literal) moves to the index position, increments the
saturated counter, and clears the flag if the counter was already $3$, while a
$0$ returns to formula level with the counter reset; at every other position it
follows the syntax automaton and leaves counter and flag unchanged. The pass never
emits.
\end{definition}

\begin{lemma}[Index bits do not count]
\lean{Complexity.satWidth_index}
\label{Complexity.satWidth_index}
\leanok
\uses{Complexity.satSyntaxStep, Complexity.satWidthStep}
\difficulty{2}
For every $n \in \mathbb{N}$, every counter value $c$, and every flag $g$, folding
the width-pass transition over the unary run $1^n$ from the index position with
counter $c$ and flag $g$ returns exactly the same state: index bits change neither
the position, nor the counter, nor the flag.
\end{lemma}

\begin{lemma}[One literal increments the counter once]
\lean{Complexity.satWidth_literal}
\label{Complexity.satWidth_literal}
\leanok
\uses{Complexity.satWidthCap, Complexity.satWidthStep, Complexity.satWidth_index, Std.Sat.CNF.serializeLit}
\difficulty{3}
Let $\ell$ be a literal, $c \in \{0,1,2,3\}$ a counter value, and $g$ a Boolean
flag. Folding the width-pass transition over the serialization of $\ell$, from the
clause position with counter $c$ and flag $g$, returns the clause position with
counter $\min(c + 1, 3)$ and flag $g \wedge [c < 3]$, independently of the
variable index and polarity of $\ell$.
\end{lemma}

\begin{lemma}[Clause width fold]
\lean{Complexity.satWidth_clause}
\label{Complexity.satWidth_clause}
\leanok
\uses{Complexity.satWidthCap, Complexity.satWidthStep, Complexity.satWidth_literal, Std.Sat.CNF.serializeClause, Std.Sat.CNF.serializeLit}
\difficulty{4}
Let $C$ be a clause, $c \in \{0,1,2,3\}$ a counter value, and $g$ a Boolean flag.
Folding the width-pass transition over the serialization of $C$, from the clause
position with counter $c$ and flag $g$, returns the formula position with counter
$0$ and flag $g \wedge [c + |C| \le 3]$, where $|C|$ counts literal occurrences
including repetitions.
\end{lemma}

\begin{lemma}[Formula width fold]
\lean{Complexity.satWidth_formula}
\label{Complexity.satWidth_formula}
\leanok
\uses{Complexity.satSyntaxStep, Complexity.satWidth, Complexity.satWidthStep, Complexity.satWidth_clause, Std.Sat.CNF.serialize, Std.Sat.CNF.serializeClause}
\difficulty{4}
Folding the width-pass transition over the serialization of a CNF formula
$\varphi$ from the initial state ends at the exact-end position with counter zero
and the goodness flag conjoined with the Boolean width-three test of $\varphi$;
every clause is checked and the empty formula passes vacuously.
\end{lemma}

\begin{definition}[Width scan]
\lean{Complexity.satWidthScan}
\label{Complexity.satWidthScan}
\leanok
\uses{Complexity.satWidthStep}
For a binary string $x$, the width scan of $x$ is the final goodness flag obtained
by folding the width-pass transition over $x$ from the initial state (formula
position, counter $0$, flag true).
\end{definition}

\begin{lemma}[Width scan on serialized formulas]
\lean{Complexity.satWidthScan_serialize}
\label{Complexity.satWidthScan_serialize}
\leanok
\uses{Complexity.satWidth, Complexity.satWidthScan, Complexity.satWidth_formula, Std.Sat.CNF.serialize}
\difficulty{1}
For every CNF formula $\varphi$, the width scan of the serialization of $\varphi$
equals the Boolean width-three test of $\varphi$, which is true exactly when every
clause of $\varphi$ has at most three literal occurrences.
\end{lemma}

\begin{lemma}[The width pass is a real machine]
\lean{Complexity.satWidthScan_poly}
\label{Complexity.satWidthScan_poly}
\leanok
\uses{Complexity.PolyTimeComputable, Complexity.satScanTM, Complexity.satScan_computes, Complexity.satWidthScan, Complexity.satWidthStep}
\difficulty{2}
The single-bit width-scan function is polynomial-time computable, realized by a
one-way finite-state scanner with no work tapes and time exactly $n + 1$.
\end{lemma}

\begin{lemma}[The 3SAT verdict is polynomial-time]
\lean{Complexity.satVerdict_true_poly}
\label{Complexity.satVerdict_true_poly}
\leanok
\uses{Complexity.PolyTimeComputable, Complexity.PolyTimeComputable.comp, Complexity.polyTimeComputable_and, Complexity.polyTimeComputable_const, Complexity.polyTimeComputable_ite, Complexity.satGood, Complexity.satInstance, Complexity.satSafeValue, Complexity.satSafeValue_poly, Complexity.satSplit, Complexity.satSplitValid, Complexity.satSyntax_spec, Complexity.satVerdict, Complexity.satWidth, Complexity.satWidthScan, Complexity.satWidthScan_poly, Complexity.satWidthScan_serialize, Complexity.satWitness, Complexity.sat_parse_repr, Complexity.sat_pipeline_poly, Std.Sat.CNF.decode, Std.Sat.CNF.decode_serialize, Std.Sat.CNF.fallback, Std.Sat.CNF.parse, Std.Sat.CNF.parse_serialize, Turing.pairDecode, Turing.pairDecode_pairEncode, Turing.solveSplit}
\difficulty{5}
The single-bit function computing the 3SAT-mode verifier verdict on a binary string
is polynomial-time computable: the machine validates the entire syntax before
running the width pass, and runs the evaluation pass only after width success;
split failure rejects, while syntax failure accepts the empty fallback.
\end{lemma}

\begin{theorem}[3SAT is in NP]
\lean{Complexity.SAT3_mem_NP}
\label{Complexity.SAT3_mem_NP}
\leanok
\uses{Complexity.NP, Complexity.SAT3, Complexity.sat3_verifier_equiv, Complexity.satVerdict_true_poly, Complexity.satVerifier, Complexity.satVerifier_of_poly}
\difficulty{2}
The language $\mathrm{3SAT}$ — binary strings whose decoded formula is a
satisfiable 3CNF — belongs to $\mathrm{NP}$, with the same certificate parameters
as $\mathrm{SAT}$: a satisfying assignment of exactly $|x| + 1$ bits, verified in
polynomial time after an additional width check. [AB09, Theorem~2.10, membership]
\end{theorem}
\begin{lemma}[Clause evaluation congruence]
\lean{Complexity.satClause_congr}
\label{Complexity.satClause_congr}
\leanok
\difficulty{2}
Let $C$ be a clause and $a, b$ two truth assignments that agree on the variable
index of every literal occurring in $C$. Then $C$ evaluates identically under $a$
and under $b$.
\end{lemma}

\begin{definition}[Clause-splitting chain]
\lean{Complexity.satChain}
\label{Complexity.satChain}
\leanok
Given a head literal $h$, a list $R$ of remaining literals, and a fresh-variable
cursor $n \in \mathbb{N}$, the clause-splitting chain is defined by recursion on
$R$ and returns a CNF formula together with the next unused cursor. If $R$ has at
least three literals $b, c, d, \dots$, the chain is the link clause
$[h,\, b,\, (n,\mathrm{true})]$ followed by the chain on head
$(n,\mathrm{false})$, remaining literals $c, d, \dots$, and cursor $n + 1$, whose
returned cursor is kept. Otherwise the chain is the single clause consisting of
$h$ followed by $R$, with the cursor $n$ unchanged.
[AB09, Section~2.3.5, proof of Lemma~2.14]
\end{definition}

\begin{definition}[Splitting one clause]
\lean{Complexity.satSplitClause}
\label{Complexity.satSplitClause}
\leanok
\uses{Complexity.satChain}
Splitting of a single clause at a fresh-variable cursor: the empty clause is
preserved as the formula containing just the empty clause (which must remain
false), and a nonempty clause is split by the chain construction.
\end{definition}

\begin{lemma}[Chain cursor monotonicity]
\lean{Complexity.satChain_cursor}
\label{Complexity.satChain_cursor}
\leanok
\uses{Complexity.satChain}
\difficulty{3}
For every head literal, remaining clause, and starting cursor $n$, the cursor
returned by the chain construction is at least $n$: the fresh-variable allocator
never moves backwards.
\end{lemma}

\begin{lemma}[Chain clauses have width three]
\lean{Complexity.satChain_width}
\label{Complexity.satChain_width}
\leanok
\uses{Complexity.satChain, Std.Sat.CNF.WidthAtMost}
\difficulty{3}
Every clause of the formula emitted by the chain construction has at most three
literals.
\end{lemma}

\begin{lemma}[Chain soundness]
\lean{Complexity.satChain_sound}
\label{Complexity.satChain_sound}
\leanok
\uses{Complexity.satChain}
\difficulty{5}
Let $a$ be a truth assignment satisfying the formula emitted by the chain
construction on a head literal, remaining clause, and cursor. Then $a$ satisfies
the original clause (head followed by the remaining literals). No freshness
hypothesis is needed in this direction.
\end{lemma}

\begin{lemma}[One splitting step extends the assignment]
\lean{Complexity.satChain_extend_step}
\label{Complexity.satChain_extend_step}
\leanok
\uses{Complexity.satClause_congr}
\difficulty{5}
Let a clause with head literal, a second literal $b$, and a tail have all its
variable indices below $n$, and let $a$ be an assignment satisfying it. Then there
is an assignment $a'$ agreeing with $a$ below $n$ such that $a'$ satisfies both the
link clause $[\mathrm{head},\, b,\, (n,\mathrm{true})]$ and the continuation clause
headed by $(n,\mathrm{false})$ followed by the tail; the fresh value $a'(n)$ is the
truth of the tail under $a$.
\end{lemma}

\begin{lemma}[Chain completeness]
\lean{Complexity.satChain_complete}
\label{Complexity.satChain_complete}
\leanok
\uses{Complexity.satChain, Complexity.satChain_extend_step, Complexity.satClause_congr}
\difficulty{5}
Let a clause (head plus remaining literals) have all variable indices below $n$,
and let $a$ be an assignment satisfying it. Then there is an assignment $a'$
agreeing with $a$ on all indices below $n$ that satisfies the formula emitted by
the chain construction at cursor $n$.
\end{lemma}

\begin{lemma}[Chain variables below the returned cursor]
\lean{Complexity.satChain_vars}
\label{Complexity.satChain_vars}
\leanok
\uses{Complexity.satChain, Complexity.satChain_cursor}
\difficulty{4}
If every variable index of the input clause (head plus remaining literals) is below
the starting cursor $n$, then every variable index occurring in the formula emitted
by the chain construction is below the returned cursor.
\end{lemma}

\begin{lemma}[Split-clause cursor monotonicity]
\lean{Complexity.satSplitClause_cursor}
\label{Complexity.satSplitClause_cursor}
\leanok
\uses{Complexity.satChain_cursor, Complexity.satSplitClause}
\difficulty{2}
For every clause $C$ and cursor $n$, the cursor returned by splitting $C$ at $n$ is
at least $n$; this includes the empty clause.
\end{lemma}

\begin{lemma}[Split clauses are 3CNF]
\lean{Complexity.satSplitClause_width}
\label{Complexity.satSplitClause_width}
\leanok
\uses{Complexity.satChain_width, Complexity.satSplitClause, Std.Sat.CNF.WidthAtMost}
\difficulty{2}
For every clause $C$ and cursor $n$, every clause of the formula obtained by
splitting $C$ at $n$ has at most three literals.
\end{lemma}

\begin{lemma}[Split-clause variable bound]
\lean{Complexity.satSplitClause_vars}
\label{Complexity.satSplitClause_vars}
\leanok
\uses{Complexity.satChain_vars, Complexity.satSplitClause}
\difficulty{2}
If every variable index of a clause $C$ is below the cursor $n$, then every
variable index occurring in the formula obtained by splitting $C$ at $n$ is below
the returned cursor.
\end{lemma}

\begin{lemma}[Split-clause soundness]
\lean{Complexity.satSplitClause_sound}
\label{Complexity.satSplitClause_sound}
\leanok
\uses{Complexity.satChain_sound, Complexity.satSplitClause}
\difficulty{2}
Every truth assignment satisfying the formula obtained by splitting a clause $C$ at
some cursor also satisfies $C$ itself.
\end{lemma}

\begin{lemma}[Split-clause completeness]
\lean{Complexity.satSplitClause_complete}
\label{Complexity.satSplitClause_complete}
\leanok
\uses{Complexity.satChain_complete, Complexity.satSplitClause}
\difficulty{2}
Let a clause $C$ have all variable indices below $n$ and let $a$ be an assignment
satisfying $C$. Then some assignment agreeing with $a$ below $n$ satisfies the
formula obtained by splitting $C$ at $n$.
\end{lemma}

\begin{definition}[Formula transform with threaded cursor]
\lean{Complexity.satTransformFrom}
\label{Complexity.satTransformFrom}
\leanok
\uses{Complexity.satSplitClause}
For a CNF formula $\varphi$ and a fresh-variable cursor $n$, the threaded clause
transform of $\varphi$ at $n$ is obtained by splitting each clause of $\varphi$ in
order, threading one global fresh-variable cursor from clause to clause starting
at $n$, and concatenating the resulting 3CNF pieces; the final cursor is returned
alongside. [AB09, Section~2.3.5, proof of Lemma~2.14]
\end{definition}

\begin{lemma}[Literal-wise bound gives the variable bound]
\lean{Complexity.sat_numVars_le}
\label{Complexity.sat_numVars_le}
\leanok
\uses{Std.Sat.CNF.numVars}
\difficulty{3}
If every variable index occurring in a CNF formula $\varphi$ is strictly below
$n$, then the variable bound of $\varphi$ (the maximum of $v + 1$ over occurring
indices $v$) is at most $n$.
\end{lemma}

\begin{lemma}[The transform is 3CNF]
\lean{Complexity.satTransformFrom_width}
\label{Complexity.satTransformFrom_width}
\leanok
\uses{Complexity.satSplitClause, Complexity.satSplitClause_width, Complexity.satTransformFrom, Std.Sat.CNF.WidthAtMost}
\difficulty{3}
For every CNF formula $\varphi$ and cursor $n$, the formula produced by the
threaded clause transform of $\varphi$ at $n$ has width at most three.
\end{lemma}

\begin{lemma}[Transform soundness]
\lean{Complexity.satTransformFrom_sound}
\label{Complexity.satTransformFrom_sound}
\leanok
\uses{Complexity.satSplitClause, Complexity.satSplitClause_sound, Complexity.satTransformFrom}
\difficulty{3}
Every truth assignment satisfying the threaded clause transform of a CNF formula
$\varphi$ also satisfies $\varphi$ itself, with the same assignment.
\end{lemma}

\begin{lemma}[Transform completeness]
\lean{Complexity.satTransformFrom_complete}
\label{Complexity.satTransformFrom_complete}
\leanok
\uses{Complexity.eval_congr_of_lt_numVars, Complexity.satSplitClause, Complexity.satSplitClause_complete, Complexity.satSplitClause_cursor, Complexity.satSplitClause_vars, Complexity.satTransformFrom, Complexity.sat_numVars_le}
\difficulty{5}
Let a CNF formula $\varphi$ have all variable indices below $n$, and let $a$ be an
assignment satisfying $\varphi$. Then some assignment agreeing with $a$ on all
indices below $n$ satisfies the threaded clause transform of $\varphi$ at cursor
$n$.
\end{lemma}

\begin{definition}[The clause-splitting transform]
\lean{Complexity.satTransform}
\label{Complexity.satTransform}
\leanok
\uses{Complexity.satTransformFrom, Std.Sat.CNF.numVars}
The formula-level SAT-to-3SAT transform: run the threaded clause transform starting
the fresh-variable cursor at exactly the variable bound of the input formula, so
that all allocated variables are genuinely fresh.
\end{definition}

\begin{lemma}[Equisatisfiability of the transform]
\lean{Complexity.satTransform_equisat}
\label{Complexity.satTransform_equisat}
\leanok
\uses{Complexity.satTransform, Complexity.satTransformFrom_complete, Complexity.satTransformFrom_sound, Complexity.sat_literal_lt_numVars, Std.Sat.CNF.Satisfiable, Std.Sat.CNF.numVars}
\difficulty{3}
For every CNF formula $\varphi$, the clause-splitting transform of $\varphi$ is
satisfiable if and only if $\varphi$ is satisfiable. [AB09, Lemma~2.14]
\end{lemma}

\begin{definition}[The string-level reduction]
\lean{Complexity.satReduction}
\label{Complexity.satReduction}
\leanok
\uses{Complexity.satTransform, Std.Sat.CNF.decode, Std.Sat.CNF.serialize}
The string-level SAT-to-3SAT reduction is the map on binary strings that decodes
its input to a CNF formula (with the empty formula as fallback for malformed
inputs), applies the clause-splitting transform (fresh variables allocated from
the decoded formula's variable bound upwards), and serializes the result.
[AB09, Lemma~2.14]
\end{definition}

\begin{lemma}[Reduction correctness]
\lean{Complexity.satReduction_correct}
\label{Complexity.satReduction_correct}
\leanok
\uses{Complexity.SAT, Complexity.SAT3, Complexity.satReduction, Complexity.satTransformFrom_width, Complexity.satTransform_equisat, Std.Sat.CNF.Satisfiable, Std.Sat.CNF.WidthAtMost, Std.Sat.CNF.decode, Std.Sat.CNF.decode_serialize, Std.Sat.CNF.numVars}
\difficulty{3}
For every binary string $x$ — with no well-formedness hypothesis — $x$ belongs to
$\mathrm{SAT}$ if and only if the image of $x$ under the string-level reduction
belongs to $\mathrm{3SAT}$.
\end{lemma}

\begin{lemma}[Reduction on malformed inputs]
\lean{Complexity.satReduction_fallback}
\label{Complexity.satReduction_fallback}
\leanok
\uses{Complexity.satReduction, Complexity.satTransform, Complexity.satTransformFrom, Std.Sat.CNF.decode, Std.Sat.CNF.fallback, Std.Sat.CNF.parse, Std.Sat.CNF.serialize}
\difficulty{1}
If the exact CNF parser fails on a binary string $x$, then the string-level
reduction maps $x$ to the serialization of the (unchanged) empty formula.
\end{lemma}

\begin{lemma}[Clause serialization length bound]
\lean{Complexity.sat_clause_measure}
\label{Complexity.sat_clause_measure}
\leanok
\uses{Std.Sat.CNF.serializeClause, Std.Sat.CNF.serializeLit}
\difficulty{2}
For every clause $C$, the serialization of $C$ has length at least $|C| + 1$,
where $|C|$ counts literal occurrences: there is room for every literal and the
terminator.
\end{lemma}
\begin{definition}[Two-tape reduction action]
\lean{Complexity.satRedAction}
\label{Complexity.satRedAction}
\leanok
\uses{Turing.Action}
Given an input move $m$, optional writes $w_0, w_1$ and head moves $d_0, d_1$ for
the two work tapes, an optional output bit $b$, and an optional next state $q$,
the two-tape reduction action is the action of the 35-state clause-splitting
transducer performing exactly these operations; work tape $0$ is a unary
fresh-variable counter and work tape $1$ a temporary literal buffer.
\end{definition}

\begin{definition}[Candidate clause-splitting transducer]
\lean{Complexity.satRedTM}
\label{Complexity.satRedTM}
\leanok
\uses{Complexity.satRedAction, Turing.FinTM}
A 35-state, two-tape candidate transducer for the clause-splitting reduction on
previously validated serialized CNF words: states 0--1 install a permanent left
marker on the literal buffer, states 2--6 silently compute the maximum unary
literal length on the counter tape, states 7--8 rewind the native input, and
states 9--34 form a proposed streaming serializer. Only the initialization, the
maximum pass, and the rewind are verified for this machine; it is not itself a
polynomial-time witness.
\end{definition}

\begin{definition}[Canonical unary counter tape]
\lean{Complexity.satRedCounter}
\label{Complexity.satRedCounter}
\leanok
\uses{Turing.FinTM.bufferTape}
For every $n \in \mathbb{N}$, the canonical unary counter tape of length $n$ is
the tape content holding a one at each offset $0, \dots, n - 1$ and a blank at
every other integer offset; its occupied length $n$ represents the next unused
variable index.
\end{definition}

\begin{definition}[Buffer with permanent marker]
\lean{Complexity.satRedBuffer}
\label{Complexity.satRedBuffer}
\leanok
\uses{Turing.FinTM.bufferTape}
The literal-buffer tape content determined by a word and a cut position: cells
before the cut have been erased, cells from the cut onwards carry the word's
remaining bits, and the cell at offset $-1$ carries a permanent marker allowing
return even after the payload has been completely erased.
\end{definition}

\begin{definition}[Canonical reduction frame]
\lean{Complexity.satRedCfg}
\label{Complexity.satRedCfg}
\leanok
\uses{Complexity.satRedCounter, Turing.Cfg}
Given an input $x$, an optional control state $q$, a position $i \le |x|$, a
counter length $n$, a buffer tape $B$, head offsets $a, b \in \mathbb{Z}$, and an
output word $o$, the canonical reduction frame is the configuration of the
35-state transducer in control $q$, with the input head just before position $i$,
the unary counter tape of length $n$ on work tape $0$ with head at $a$, the buffer
tape $B$ on work tape $1$ with head at $b$, and accumulated output $o$.
\end{definition}

\begin{lemma}[Input symbol of the reduction frame]
\lean{Complexity.satRedCfg_input}
\label{Complexity.satRedCfg_input}
\leanok
\uses{Complexity.satRedCfg, Turing.Cfg.inputSymbol, Turing.FinTM.inputSymbol_at}
\difficulty{1}
Let $x$ be an input and $i \le |x|$. In every canonical reduction frame on input
$x$ at input position $i$, the symbol read by the input head is the $i$-th bit of
$x$ when $i < |x|$, and the blank when $i = |x|$.
\end{lemma}

\begin{lemma}[Reading the unary counter]
\lean{Complexity.satRedCounter_read}
\label{Complexity.satRedCounter_read}
\leanok
\uses{Complexity.satRedCounter, Turing.FinTM.bufferTape_nat}
\difficulty{1}
The unary counter tape of length $n$ holds a one at every nonnegative offset
$j < n$ and a blank at every offset $j \ge n$: its occupied cells are exactly its
unary prefix.
\end{lemma}

\begin{lemma}[Counter left boundary]
\lean{Complexity.satRedCounter_left}
\label{Complexity.satRedCounter_left}
\leanok
\uses{Complexity.satRedCounter, Turing.FinTM.bufferTape_left}
\difficulty{1}
The unary counter tape of any length is blank at offset $-1$.
\end{lemma}

\begin{lemma}[Writing the counter]
\lean{Complexity.satRedCounter_write}
\label{Complexity.satRedCounter_write}
\leanok
\uses{Complexity.satRedCounter, Complexity.satRedCounter_read, Turing.FinTM.bufferTape_append}
\difficulty{3}
Let $n, j \in \mathbb{N}$ with $j \le n$. Writing a one at offset $j$ of the unary
counter tape of length $n$ yields the unary counter tape of length
$\max(n, j + 1)$: writing within the current prefix preserves it, while writing
its right blank extends the length by one.
\end{lemma}

\begin{lemma}[Action calculus on reduction frames]
\lean{Complexity.satRedAction_apply}
\label{Complexity.satRedAction_apply}
\leanok
\uses{Complexity.satRedAction, Complexity.satRedCfg, Complexity.satRedCounter, Turing.Action.apply, Turing.moveInputPos}
\difficulty{4}
Applying a two-tape reduction action to a canonical reduction frame yields the
canonical frame whose input position, counter tape, buffer tape, both head
offsets, and appended output are exactly as prescribed by the action's components;
this frame-level calculus is shared by the maximum pass and the streaming
serializer.
\end{lemma}

\begin{lemma}[Silent move of the transducer]
\lean{Complexity.satRed_move}
\label{Complexity.satRed_move}
\leanok
\uses{Complexity.satRedAction, Complexity.satRedAction_apply, Complexity.satRedCfg, Complexity.satRedCfg_input, Complexity.satRedCounter, Complexity.satRedTM, Turing.Action.apply, Turing.Cfg.workTapeSymbols, Turing.MultiTapeTM.step, Turing.moveInputPos}
\difficulty{4}
If the transducer's transition at the current state, input symbol, and work
symbols is a silent action with no writes, then one machine step from the
canonical reduction frame changes only the control state and the three head
positions, as prescribed; the counter length, buffer, and output are unchanged.
\end{lemma}

\begin{lemma}[One step equals a one-step run]
\lean{Complexity.satRed_one}
\label{Complexity.satRed_one}
\leanok
\uses{Complexity.satRedTM, Turing.Cfg, Turing.MultiTapeTM.runFrom, Turing.MultiTapeTM.runFrom_succ_eq_step, Turing.MultiTapeTM.runFrom_zero, Turing.MultiTapeTM.step}
\difficulty{1}
Running the clause-splitting transducer for exactly one step from any
configuration coincides with applying its single-step transition.
\end{lemma}

\begin{lemma}[Counter rewind of the transducer]
\lean{Complexity.satRed_counterBack}
\label{Complexity.satRed_counterBack}
\leanok
\uses{Complexity.satRedAction, Complexity.satRedCfg, Complexity.satRedCounter_left, Complexity.satRedCounter_read, Complexity.satRedTM, Complexity.satRed_move, Turing.MultiTapeTM.runFrom, Turing.MultiTapeTM.runFrom_succ_eq_step, Turing.MultiTapeTM.runFrom_zero, Turing.MultiTapeTM.step, Turing.moveInputPos_zero}
\difficulty{5}
Let $q$ be a transducer state whose transition moves the counter head left while
its cell is occupied and otherwise steps right into a destination state. For every
$r \le n$, starting from the canonical frame with counter length $n$ and counter
head at offset $r - 1$, exactly $r + 1$ silent transitions reach the destination
state with the counter head at $0$; input position, buffer, and output are
untouched.
\end{lemma}

\begin{lemma}[Maximum pass over a unary run]
\lean{Complexity.satRed_maxOnes}
\label{Complexity.satRed_maxOnes}
\leanok
\uses{Complexity.satRedAction, Complexity.satRedAction_apply, Complexity.satRedCfg, Complexity.satRedCfg_input, Complexity.satRedCounter_write, Complexity.satRedTM, Turing.Action.apply, Turing.MultiTapeTM.runFrom, Turing.MultiTapeTM.runFrom_succ_eq_step, Turing.MultiTapeTM.runFrom_zero, Turing.MultiTapeTM.step, Turing.moveInputPos_pos_of_ne_right}
\difficulty{5}
Suppose the input contains a unary run $1^v$ at the current position and the
transducer is in its index-scanning state with counter length $n$ and counter head
at $j \le n$. Then exactly $v$ steps consume the run, advancing the input by $v$,
moving the counter head to $j + v$, and extending the counter length to
$\max(n,\, j + v)$.
\end{lemma}

\begin{lemma}[Maximum pass over one literal]
\lean{Complexity.satRed_maxLiteral}
\label{Complexity.satRed_maxLiteral}
\leanok
\uses{Complexity.satRedAction, Complexity.satRedAction_apply, Complexity.satRedCfg, Complexity.satRedCfg_input, Complexity.satRedCounter_write, Complexity.satRedTM, Complexity.satRed_counterBack, Complexity.satRed_maxOnes, Complexity.satRed_move, Complexity.satRed_one, Std.Sat.CNF.serializeLit, Turing.Action.apply, Turing.MultiTapeTM.runFrom, Turing.MultiTapeTM.runFrom_add, Turing.MultiTapeTM.step, Turing.moveInputPos_pos_of_ne_right}
\difficulty{6}
Suppose the input contains a serialized literal $(v,\mathrm{pol})$ at the current
position and the transducer is in its clause-level maximum state with counter
length $n$ and counter head at $0$. Then exactly $2v + 5$ transitions consume the
literal, record the counter length $\max(n,\, v + 1)$, restore the counter head to
$0$, and return to the same state with the input advanced by $v + 3$.
\end{lemma}

\begin{lemma}[Maximum folding over concatenation]
\lean{Complexity.sat_foldMax_append}
\label{Complexity.sat_foldMax_append}
\leanok
\difficulty{2}
For lists of natural numbers $a$ and $b$, folding the maximum (with base $0$) over
the concatenation $a \mathbin{+\!\!+} b$ equals the maximum of the two separate
folds.
\end{lemma}

\begin{definition}[Clause variable maximum]
\lean{Complexity.satClauseVars}
\label{Complexity.satClauseVars}
\leanok
The maximum, over all literals of a clause, of the literal's variable index plus
one (with base $0$); this is the clause-level accumulator of the scanner
invariant.
\end{definition}

\begin{lemma}[Variable bound of a cons formula]
\lean{Complexity.sat_numVars_cons}
\label{Complexity.sat_numVars_cons}
\leanok
\uses{Complexity.satClauseVars, Complexity.sat_foldMax_append, Std.Sat.CNF.numVars}
\difficulty{1}
The variable bound of a formula with first clause $C$ and remaining clauses
$\varphi$ equals the maximum of the clause variable maximum of $C$ with the
variable bound of $\varphi$.
\end{lemma}

\begin{lemma}[Maximum pass over one clause]
\lean{Complexity.satRed_maxClause}
\label{Complexity.satRed_maxClause}
\leanok
\uses{Complexity.satClauseVars, Complexity.satRedCfg, Complexity.satRedTM, Complexity.satRed_maxLiteral, Complexity.satRed_move, Std.Sat.CNF.serializeClause, Std.Sat.CNF.serializeLit, Turing.MultiTapeTM.runFrom, Turing.MultiTapeTM.runFrom_add, Turing.MultiTapeTM.runFrom_succ_eq_step, Turing.MultiTapeTM.runFrom_zero, Turing.moveInputPos_pos_of_ne_right}
\difficulty{5}
Let $C$ be a clause with serialization $\mathrm{ser}(C)$, and let $m_C$ be the
maximum of $v + 1$ over the literals $(v, b)$ of $C$ (or $0$ if $C$ is empty).
Suppose the input is $x = p \mathbin{+\!\!+} \mathrm{ser}(C) \mathbin{+\!\!+} r$,
and the transducer is in its clause-level maximum state at input position $|p|$,
with counter length $n$, counter head at $0$, and arbitrary buffer, buffer head,
and output. Then there is a number $t \le 2\,|\mathrm{ser}(C)|$ of silent steps
after which the transducer is in formula-level maximum control at input position
$|p| + |\mathrm{ser}(C)|$, with counter length $\max(n, m_C)$, counter head back
at $0$, and buffer, buffer head, and output unchanged.
\end{lemma}

\begin{lemma}[Maximum pass over a formula]
\lean{Complexity.satRed_maxFormula}
\label{Complexity.satRed_maxFormula}
\leanok
\uses{Complexity.satClauseVars, Complexity.satRedCfg, Complexity.satRedTM, Complexity.satRed_maxClause, Complexity.satRed_move, Complexity.satRed_one, Complexity.sat_numVars_cons, Std.Sat.CNF.numVars, Std.Sat.CNF.serialize, Std.Sat.CNF.serializeClause, Turing.MultiTapeTM.runFrom, Turing.MultiTapeTM.runFrom_add, Turing.MultiTapeTM.runFrom_succ_eq_step, Turing.MultiTapeTM.runFrom_zero, Turing.MultiTapeTM.step, Turing.moveInputPos_pos_of_ne_right, Turing.moveInputPos_zero}
\difficulty{5}
Let $\varphi$ be a CNF formula with serialization $\mathrm{ser}(\varphi)$ and
its variable bound $N(\varphi)$ (one more than the largest variable index occurring in $\varphi$, or $0$ if no variable occurs). Suppose the input is
$x = p \mathbin{+\!\!+} \mathrm{ser}(\varphi) \mathbin{+\!\!+} r$, and the
transducer is in formula-level maximum control at input position $|p|$, with
counter length $n$, counter head at $0$, and arbitrary buffer, buffer head, and
output. Then there is a number $t \le 2\,|\mathrm{ser}(\varphi)|$ of silent steps
after which the transducer is at its rewind seam at input position
$|p| + |\mathrm{ser}(\varphi)| - 1$, with counter length $\max(n, N(\varphi))$,
counter head at $0$, and buffer, buffer head, and output unchanged.
\end{lemma}

\begin{lemma}[The empty buffer]
\lean{Complexity.satRedBuffer_empty}
\label{Complexity.satRedBuffer_empty}
\leanok
\uses{Complexity.satRedBuffer, Turing.FinTM.bufferTape}
\difficulty{2}
The literal buffer holding the empty word with cut $0$ is the all-blank tape
updated with the permanent marker at offset $-1$.
\end{lemma}

\begin{lemma}[Marker installation]
\lean{Complexity.satRed_init}
\label{Complexity.satRed_init}
\leanok
\uses{Complexity.satRedAction, Complexity.satRedAction_apply, Complexity.satRedBuffer, Complexity.satRedBuffer_empty, Complexity.satRedCfg, Complexity.satRedCounter, Complexity.satRedTM, Complexity.satRed_move, Complexity.satRed_one, Turing.Action.apply, Turing.FinTM.bufferTape, Turing.MultiTapeTM.initCfg, Turing.MultiTapeTM.runFrom, Turing.MultiTapeTM.runFrom_succ_eq_step, Turing.MultiTapeTM.step, Turing.moveInputPos_zero}
\difficulty{4}
For every input $x$, the first two steps of the clause-splitting transducer from
its initial configuration on $x$ are silent and yield the canonical frame in the
maximum-pass entry state with input head at $0$, empty counter, the empty marked
buffer (blank except for the permanent marker at offset $-1$), both work heads at
$0$, and empty output.
\end{lemma}

\begin{lemma}[Transducer startup seam]
\lean{Complexity.satRed_start}
\label{Complexity.satRed_start}
\leanok
\uses{Complexity.satRedAction, Complexity.satRedBuffer, Complexity.satRedCfg, Complexity.satRedTM, Complexity.satRed_init, Complexity.satRed_maxFormula, Std.Sat.CNF.numVars, Std.Sat.CNF.serialize, Turing.FinTM.controlAction, Turing.FinTM.timed_rewind, Turing.MultiTapeTM.initCfg, Turing.MultiTapeTM.runFrom, Turing.MultiTapeTM.runFrom_add}
\difficulty{5}
Let $\varphi$ be a CNF formula with serialization $\mathrm{ser}(\varphi)$ and
its variable bound $N(\varphi)$ (one more than the largest variable index occurring in $\varphi$, or $0$ if no variable occurs). There is a time $t \le 3\,|\mathrm{ser}(\varphi)| + 5$ at which the
clause-splitting transducer, started on $\mathrm{ser}(\varphi)$, is in its
streaming entry state with exactly $1^{N(\varphi)}$ on the counter tape, the
buffer empty apart from its permanent marker, the input head back at $0$, both
work heads at $0$, and empty output.
\end{lemma}
\begin{definition}[Streaming seam state]
\lean{Complexity.SatStreamState}
\label{Complexity.SatStreamState}
\leanok
A streaming seam state is a triple consisting of a fresh-variable cursor in
$\mathbb{N}$, a number of consumed input bits in $\mathbb{N}$, and one of five
grammar phases (formula, first literal, second literal, tail, finished); it is the
persistent state carried across one round of the normalized streaming schedule.
\end{definition}

\begin{lemma}[Literal parser round trip]
\lean{Complexity.satStream_parseLit}
\label{Complexity.satStream_parseLit}
\leanok
\uses{Std.Sat.CNF.parseLit, Std.Sat.CNF.serializeLit, Std.Sat.CNF.takeTrues}
\difficulty{3}
For every literal $\ell$ and remainder word $r$, the literal parser applied to the
serialization of $\ell$ followed by $r$ succeeds and returns exactly $(\ell, r)$.
\end{lemma}

\begin{definition}[One normalized streaming round]
\lean{Complexity.satStreamRound}
\label{Complexity.satStreamRound}
\leanok
\uses{Complexity.SatStreamState, Std.Sat.CNF.parseLit, Std.Sat.CNF.serializeLit}
One round of the streaming clause-splitting schedule on an input string and a seam
state: the finished phase is absorbing with empty output; an exhausted input emits
the fallback terminator $0$ and finishes; a terminator bit $0$ is copied, closing
a clause or the whole formula; a clause marker $1$ at formula phase is copied; and
otherwise a complete literal is parsed together with a one-bit lookahead — in the
tail phase with more literals pending, the round allocates one fresh variable and
emits the positive link, two clause markers, the fresh negative literal, and the
consumed original literal, while in all other literal positions the literal is
copied verbatim. Local parse failures terminate silently.
\end{definition}

\begin{definition}[Iterated rounds]
\lean{Complexity.satStreamRun}
\label{Complexity.satStreamRun}
\leanok
\uses{Complexity.SatStreamState, Complexity.satStreamRound}
Execution of a prescribed number of normalized streaming rounds from a seam state,
returning the final state together with the concatenation of all emitted chunks in
emission order. This is a pure schedule, not a time-computability claim.
\end{definition}

\begin{lemma}[Splitting the round count]
\lean{Complexity.satStreamRun_add}
\label{Complexity.satStreamRun_add}
\leanok
\uses{Complexity.SatStreamState, Complexity.satStreamRun}
\difficulty{3}
Running $m + n$ streaming rounds from a state equals running $m$ rounds, then $n$
further rounds from the intermediate state, with the two output words
concatenated.
\end{lemma}

\begin{lemma}[Finished rounds are empty]
\lean{Complexity.satStreamRun_finished}
\label{Complexity.satStreamRun_finished}
\leanok
\uses{Complexity.satStreamRound, Complexity.satStreamRun}
\difficulty{2}
For every input $x$, every $k \in \mathbb{N}$, and every seam state in the
finished phase, running $k$ streaming rounds returns the same state and the empty
output: padding after the formula terminator contributes no further output.
\end{lemma}

\begin{definition}[Pending tail output]
\lean{Complexity.satStreamTail}
\label{Complexity.satStreamTail}
\leanok
\uses{Std.Sat.CNF.serializeLit}
For a remaining clause tail and a fresh cursor $j$, the pending tail output is the
output still to be emitted after the first two literals of a clause, together with
the final cursor. An empty tail emits just the clause terminator $0$; a single
literal is copied and followed by the terminator; otherwise one link is closed
with the fresh positive literal $(j,\mathrm{true})$, a new clause is opened with
$(j,\mathrm{false})$, the next original literal is copied, and the recursion
continues on the rest of the tail at cursor $j + 1$.
[AB09, Section~2.3.5, proof of Lemma~2.14]
\end{definition}

\begin{definition}[Emitted clause body]
\lean{Complexity.satStreamClause}
\label{Complexity.satStreamClause}
\leanok
\uses{Complexity.satStreamTail, Std.Sat.CNF.serializeLit}
For a clause $C$ and fresh cursor $j$: the part of the transformed clause's
serialization emitted after the initial clause marker (which is handled by the
formula phase), including its final terminator, together with the final cursor.
Up to two leading literals are copied verbatim before the pending tail output
takes over.
\end{definition}

\begin{lemma}[Tail output matches the chain]
\lean{Complexity.satStreamTail_chain}
\label{Complexity.satStreamTail_chain}
\leanok
\uses{Complexity.satChain, Complexity.satStreamTail, Std.Sat.CNF.serializeClause, Std.Sat.CNF.serializeLit}
\difficulty{4}
For literals $a, b$, a clause tail $C$, and a cursor $j$: concatenating a clause
marker with the serialization of each clause of the chain construction on $a$,
$b :: C$, $j$ yields exactly one marker followed by the serializations of $a$ and
$b$ and the pending tail output of $C$ at $j$; moreover the chain's returned
cursor equals the tail's final cursor.
\end{lemma}

\begin{lemma}[Clause schedule realizes clause splitting]
\lean{Complexity.satStreamClause_split}
\label{Complexity.satStreamClause_split}
\leanok
\uses{Complexity.satChain, Complexity.satSplitClause, Complexity.satStreamClause, Complexity.satStreamTail_chain, Std.Sat.CNF.serializeClause}
\difficulty{2}
For every clause $C$ and cursor $j$: concatenating a clause marker with the
serialization of each clause of the split of $C$ at $j$ yields exactly one marker
followed by the emitted clause body of $C$ at $j$, and the split's returned cursor
equals the body's final cursor. This includes the empty clause and all clauses of
width at most three.
\end{lemma}

\begin{lemma}[Round at a terminator]
\lean{Complexity.satStreamRound_false}
\label{Complexity.satStreamRound_false}
\leanok
\uses{Complexity.satStreamRound}
\difficulty{1}
Let the input be $p \mathbin{+\!\!+} 0 :: r$ and the seam state have consumed
$|p|$ bits in a non-finished phase $q$. The round emits $[0]$, consumes one bit,
and moves to the finished phase when $q$ was the formula phase and back to the
formula phase otherwise.
\end{lemma}

\begin{lemma}[Round at a clause marker]
\lean{Complexity.satStreamRound_marker}
\label{Complexity.satStreamRound_marker}
\leanok
\uses{Complexity.satStreamRound}
\difficulty{1}
Let the input be $p \mathbin{+\!\!+} 1 :: r$ and the seam state have consumed
$|p|$ bits in the formula phase. The round emits $[1]$, consumes one bit, and
moves to the first-literal phase.
\end{lemma}

\begin{lemma}[Round at a complete literal]
\lean{Complexity.satStreamRound_literal}
\label{Complexity.satStreamRound_literal}
\leanok
\uses{Complexity.satStreamRound, Complexity.satStream_parseLit, Std.Sat.CNF.serializeLit}
\difficulty{4}
Let the input be $p$ followed by a serialized literal $\ell$ and a remainder $r$,
and let the seam state have consumed $|p|$ bits in a literal phase (neither
formula nor finished). If the phase is the tail phase and $r$ starts with a clause
bit $1$, the round allocates the fresh cursor $j$: it emits the positive link
$(j,\mathrm{true})$, two markers, the negative literal $(j,\mathrm{false})$, and
$\ell$, incrementing the cursor; otherwise it emits $\ell$ verbatim and advances
the phase. In both cases the consumed length advances by the serialized literal
length, so no literal buffer survives a round boundary.
\end{lemma}

\begin{lemma}[Tail phase schedule]
\lean{Complexity.satStreamRun_tail}
\label{Complexity.satStreamRun_tail}
\leanok
\uses{Complexity.satStreamRound_false, Complexity.satStreamRound_literal, Complexity.satStreamRun, Complexity.satStreamTail, Std.Sat.CNF.serializeClause, Std.Sat.CNF.serializeLit}
\difficulty{5}
Let the input contain a serialized clause tail $C$ at the consumed position, with
the seam in the tail phase at cursor $j$. Then exactly $|C| + 1$ rounds — one per
remaining literal plus one for the clause terminator — consume the serialized
tail, emit exactly the pending tail output of $C$ at $j$, return to the formula
phase, and leave the cursor at the tail's final value.
\end{lemma}

\begin{lemma}[Clause body schedule]
\lean{Complexity.satStreamRun_clause}
\label{Complexity.satStreamRun_clause}
\leanok
\uses{Complexity.satStreamClause, Complexity.satStreamRound_false, Complexity.satStreamRound_literal, Complexity.satStreamRun, Complexity.satStreamRun_tail, Std.Sat.CNF.serializeClause, Std.Sat.CNF.serializeLit}
\difficulty{5}
Let the input contain a serialized clause $C$ at the consumed position, with the
seam in the first-literal phase at cursor $j$. Then exactly $|C| + 1$ rounds
consume the serialized clause, emit exactly the emitted clause body of $C$ at
$j$, return to the formula phase, and leave the cursor at the body's final value.
\end{lemma}

\begin{definition}[Round count of a formula]
\lean{Complexity.satStreamCount}
\label{Complexity.satStreamCount}
\leanok
For a CNF formula $\varphi$ with clauses $C_1, \dots, C_m$, the round count of
$\varphi$ is $1 + \sum_{k=1}^{m} (|C_k| + 2)$: per clause one marker round, one
round per literal occurrence, and one terminator round, plus one final
formula-terminator round.
\end{definition}

\begin{lemma}[Formula schedule emits the transform]
\lean{Complexity.satStreamRun_formula}
\label{Complexity.satStreamRun_formula}
\leanok
\uses{Complexity.satStreamClause, Complexity.satStreamClause_split, Complexity.satStreamCount, Complexity.satStreamRound_false, Complexity.satStreamRound_marker, Complexity.satStreamRun, Complexity.satStreamRun_add, Complexity.satStreamRun_clause, Complexity.satTransformFrom, Std.Sat.CNF.serialize, Std.Sat.CNF.serializeClause}
\difficulty{5}
Let the input contain the serialization of a formula $\varphi$ at the consumed
position, with the seam in the formula phase at cursor $j$. Then exactly the
round count of $\varphi$ many rounds consume the serialized formula, emit exactly
the serialization of the threaded clause transform of $\varphi$ at $j$, reach the
finished phase, and leave the cursor at the transform's final value.
\end{lemma}

\begin{lemma}[Round count bound]
\lean{Complexity.satStreamCount_le}
\label{Complexity.satStreamCount_le}
\leanok
\uses{Complexity.satStreamCount, Complexity.sat_clause_measure, Std.Sat.CNF.serialize, Std.Sat.CNF.serializeClause}
\difficulty{3}
For every CNF formula $\varphi$, the round count of $\varphi$ is at most the
length of the serialization of $\varphi$: every non-finished round consumes at
least one input bit, including for empty clauses and the empty formula.
\end{lemma}

\begin{lemma}[Serialized schedule with padding]
\lean{Complexity.satStreamRun_serialize}
\label{Complexity.satStreamRun_serialize}
\leanok
\uses{Complexity.satStreamCount, Complexity.satStreamCount_le, Complexity.satStreamRun, Complexity.satStreamRun_add, Complexity.satStreamRun_finished, Complexity.satStreamRun_formula, Complexity.satTransform, Std.Sat.CNF.numVars, Std.Sat.CNF.serialize}
\difficulty{4}
Let $\varphi$ be a CNF formula with serialization $\mathrm{ser}(\varphi)$ and
its variable bound $N(\varphi)$ (one more than the largest variable index occurring in $\varphi$, or $0$ if no variable occurs). Running $|\mathrm{ser}(\varphi)| + 1$ streaming rounds on input
$\mathrm{ser}(\varphi)$, starting from the formula phase with fresh cursor
$N(\varphi)$ and nothing consumed, emits a word equal to the serialization of the
clause-splitting transform of $\varphi$: the semantic schedule finishes within the
budget and all remaining rounds are empty.
\end{lemma}

\begin{lemma}[Round output as an indexed concatenation]
\lean{Complexity.satStreamRun_output}
\label{Complexity.satStreamRun_output}
\leanok
\uses{Complexity.SatStreamState, Complexity.satStreamRound, Complexity.satStreamRun}
\difficulty{3}
For every input $x$, round count $k$, and seam state $s$, the output of $k$
streaming rounds equals the concatenation over $i < k$ of the chunk emitted by
the round applied to the $i$-fold iterate of the next-state map on $s$ — the
range-indexed form required by the emitting-loop contract.
\end{lemma}

\begin{definition}[Encoded seam word]
\lean{Complexity.satStreamWord}
\label{Complexity.satStreamWord}
\leanok
\uses{Complexity.SatStreamState, Turing.pairEncode}
The pair-encoded word representing a seam state: the unary fresh cursor paired
with the pair of the unary consumed length and a bounded unary phase tag. It
contains no round-local literal buffer.
\end{definition}

\begin{definition}[First pair projection]
\lean{Complexity.satStreamFst}
\label{Complexity.satStreamFst}
\leanok
\uses{Turing.pairDecode}
For a word $w$, the first pair projection of $w$ is the first component of $w$
decoded as a pair encoding, and is the empty word when $w$ is not a valid pair
encoding.
\end{definition}

\begin{definition}[Second pair projection]
\lean{Complexity.satStreamSnd}
\label{Complexity.satStreamSnd}
\leanok
\uses{Turing.pairDecode}
For a word $w$, the second pair projection of $w$ is the second component of $w$
decoded as a pair encoding, and is the empty word when $w$ is not a valid pair
encoding.
\end{definition}

\begin{definition}[Decoding a seam word]
\lean{Complexity.satStreamRead}
\label{Complexity.satStreamRead}
\leanok
\uses{Complexity.SatStreamState, Complexity.satStreamFst, Complexity.satStreamSnd}
The total decoding of a word into a seam state: the lengths of the two unary
counter fields give the fresh cursor and consumed length, and the length of the
phase field is reduced modulo five. The modulus only totalizes malformed words and
is the identity on all reachable tags.
\end{definition}

\begin{lemma}[Seam encoding round trip]
\lean{Complexity.satStreamRead_word}
\label{Complexity.satStreamRead_word}
\leanok
\uses{Complexity.SatStreamState, Complexity.satStreamFst, Complexity.satStreamRead, Complexity.satStreamSnd, Complexity.satStreamWord, Turing.pairDecode_pairEncode}
\difficulty{2}
For every seam state $s$, decoding the encoded seam word of $s$ returns exactly
$s$: the fresh cursor, the consumed length, and the phase are all recovered.
\end{lemma}

\begin{lemma}[Length of a seam word]
\lean{Complexity.satStreamWord_length}
\label{Complexity.satStreamWord_length}
\leanok
\uses{Complexity.SatStreamState, Complexity.satStreamWord, Turing.pairEncode}
\difficulty{2}
For every seam state with fresh cursor $f$, consumed length $c$, and phase
$q \in \{0, \dots, 4\}$, the encoded seam word has length exactly
$2f + 2c + q + 4$; in particular it is linear in the two unary counters.
\end{lemma}

\begin{definition}[Length invariant of the loop]
\lean{Complexity.satStreamBound}
\label{Complexity.satStreamBound}
\leanok
\uses{Complexity.SatStreamState}
The length-only loop invariant on an input $x$ and a seam state: the consumed
length is at most $|x|$, and the fresh cursor is at most $|x|$ plus the consumed
length — at most one fresh variable is allocated per consumed bit.
\end{definition}

\begin{lemma}[Rounds preserve the invariant]
\lean{Complexity.satStreamRound_bound}
\label{Complexity.satStreamRound_bound}
\leanok
\uses{Complexity.SatStreamState, Complexity.satStreamBound, Complexity.satStreamRound, Complexity.sat_parseLit_repr, Std.Sat.CNF.parseLit, Std.Sat.CNF.serializeLit}
\difficulty{4}
If a seam state satisfies the length invariant on an input $x$, then so does the
state after one streaming round — including on malformed local inputs, where
finished and failure cases preserve the counters.
\end{lemma}

\begin{lemma}[Size of reachable seam words]
\lean{Complexity.satStreamBound_size}
\label{Complexity.satStreamBound_size}
\leanok
\uses{Complexity.SatStreamState, Complexity.satStreamBound, Complexity.satStreamWord, Complexity.satStreamWord_length}
\difficulty{2}
If a seam state satisfies the length invariant on an input $x$, then its fresh
cursor is at most $2\,|x|$ and its encoded seam word has length at most
$6\,|x| + 8$.
\end{lemma}

\begin{definition}[Startup seam state]
\lean{Complexity.satStreamStart}
\label{Complexity.satStreamStart}
\leanok
\uses{Complexity.SatStreamState, Complexity.satSyntax, Std.Sat.CNF.decode, Std.Sat.CNF.numVars}
The seam state selected before any emission: if the complete syntax scan of the
input accepts, start in the formula phase with the fresh cursor at the decoded
formula's variable bound and nothing consumed; otherwise start in the formula
phase at the right input boundary, whose first round emits only the fallback
terminator and then finishes — even at input length zero.
\end{definition}

\begin{lemma}[Startup satisfies the invariant]
\lean{Complexity.satStreamStart_bound}
\label{Complexity.satStreamStart_bound}
\leanok
\uses{Complexity.satStreamBound, Complexity.satStreamStart, Std.Sat.CNF.numVars_decode_le}
\difficulty{2}
For every input string $x$, the startup seam state of $x$ satisfies the length
invariant on $x$; this holds both for the valid and for the fallback startup
state.
\end{lemma}

\begin{lemma}[Fallback output]
\lean{Complexity.satStreamRun_fallback}
\label{Complexity.satStreamRun_fallback}
\leanok
\uses{Complexity.satStreamRound, Complexity.satStreamRun, Complexity.satStreamRun_finished}
\difficulty{1}
Any positive number of streaming rounds from the right-boundary fallback state
emits exactly the single bit $0$: no prefix of the malformed string is emitted,
including on trailing-data failures.
\end{lemma}

\begin{lemma}[The schedule computes the reduction]
\lean{Complexity.satStreamRun_correct}
\label{Complexity.satStreamRun_correct}
\leanok
\uses{Complexity.satReduction, Complexity.satReduction_fallback, Complexity.satStreamRun, Complexity.satStreamRun_fallback, Complexity.satStreamRun_serialize, Complexity.satStreamStart, Complexity.satSyntax, Complexity.satSyntax_spec, Complexity.sat_parse_repr, Std.Sat.CNF.decode, Std.Sat.CNF.parse}
\difficulty{4}
For every binary string $x$, running $|x| + 1$ streaming rounds from the startup
seam state of $x$ emits exactly the image of $x$ under the string-level
SAT-to-3SAT reduction. This is an exact output identity on every string, prior to
any machine-computability claim.
\end{lemma}

\begin{definition}[Encoded next-state function]
\lean{Complexity.satStreamStep}
\label{Complexity.satStreamStep}
\leanok
\uses{Complexity.satStreamRead, Complexity.satStreamRound, Complexity.satStreamWord}
For an input $x$ and a word $w$, the encoded next-state function returns the
encoded seam word of the state obtained by decoding $w$ and applying one streaming
round on input $x$.
\end{definition}

\begin{definition}[Encoded chunk function]
\lean{Complexity.satStreamEmit}
\label{Complexity.satStreamEmit}
\leanok
\uses{Complexity.satStreamRead, Complexity.satStreamRound}
The word-level chunk function consumed by the audited emitter interface: decode a
seam word, apply one streaming round, and return the emitted chunk.
\end{definition}

\begin{definition}[Encoded loop invariant]
\lean{Complexity.satStreamInv}
\label{Complexity.satStreamInv}
\leanok
\uses{Complexity.satStreamBound, Complexity.satStreamWord}
The loop invariant on an input $x$ and a word $w$: $w$ is the canonical encoding
of some seam state satisfying the length invariant on $x$. Arbitrary malformed
state words are not admitted as seams.
\end{definition}

\begin{lemma}[Encoded invariant at startup]
\lean{Complexity.satStreamInv_start}
\label{Complexity.satStreamInv_start}
\leanok
\uses{Complexity.satStreamInv, Complexity.satStreamStart, Complexity.satStreamStart_bound, Complexity.satStreamWord}
\difficulty{1}
For every input string $x$, the encoded startup seam word of $x$ satisfies the
encoded loop invariant.
\end{lemma}

\begin{lemma}[Encoded invariant is preserved]
\lean{Complexity.satStreamInv_step}
\label{Complexity.satStreamInv_step}
\leanok
\uses{Complexity.satStreamInv, Complexity.satStreamRead_word, Complexity.satStreamRound, Complexity.satStreamRound_bound, Complexity.satStreamStep}
\difficulty{2}
If a word satisfies the encoded loop invariant on an input $x$, then so does its
image under the encoded next-state function.
\end{lemma}

\begin{lemma}[Encoded and structured orbits agree]
\lean{Complexity.satStreamStep_iterate}
\label{Complexity.satStreamStep_iterate}
\leanok
\uses{Complexity.SatStreamState, Complexity.satStreamRead_word, Complexity.satStreamRound, Complexity.satStreamStep, Complexity.satStreamWord}
\difficulty{3}
For every input $x$, seam state $s$, and $k$, the $k$-fold iterate of the encoded
next-state function on the encoded word of $s$ equals the encoded word of the
$k$-fold iterate of the structured round map on $s$.
\end{lemma}

\begin{lemma}[Output identity for the emitting loop]
\lean{Complexity.satStream_output_identity}
\label{Complexity.satStream_output_identity}
\leanok
\uses{Complexity.satReduction, Complexity.satStreamEmit, Complexity.satStreamRead_word, Complexity.satStreamRun_correct, Complexity.satStreamRun_output, Complexity.satStreamStart, Complexity.satStreamStep, Complexity.satStreamStep_iterate, Complexity.satStreamWord}
\difficulty{2}
For every binary string $x$, concatenating over $i < |x| + 1$ the chunks produced
by the encoded chunk function at the $i$-fold iterate of the encoded next-state
function on the encoded startup word equals the image of $x$ under the
string-level SAT-to-3SAT reduction — exactly the expression returned by the
emitting-loop constructor with round budget $R(n) = n$.
\end{lemma}
\begin{definition}[Pairing administrative action]
\lean{Complexity.satPairAction}
\label{Complexity.satPairAction}
\leanok
\uses{Turing.Action, Turing.FinTM}
Let $M$ be a base machine with $k$ work tapes. Given an input move $m$, a
capture-head move $d$, an optional output bit $b$, and an optional target state
$q$, the pairing administrative action is the action of the pairing machine over
$M$ that writes nothing, keeps the completed source bank (the first $k$ work
heads) stationary, moves the input head by $m$ and the last (capture) head by $d$,
emits $b$ when present, and switches to $q$.
\end{definition}

\begin{definition}[Input-retaining pairing machine]
\lean{Complexity.satPairTM}
\label{Complexity.satPairTM}
\leanok
\uses{Complexity.satPairAction, Turing.FinTM, Turing.FinTM.controlAction, Turing.captureAction}
Given a machine $M$ computing a function $f$, the pairing machine with one extra
tape: it runs $M$, capturing its output onto the fresh last tape, rewinds the
native input and the captured word, replays each captured bit twice, emits the
pairing separator $01$, and finally copies the native input. On input $x$ its
output is therefore the pair encoding of $f(x)$ with $x$.
\end{definition}

\begin{definition}[Pairing configuration]
\lean{Complexity.satPairCfg}
\label{Complexity.satPairCfg}
\leanok
\uses{Complexity.satPairTM, Turing.Cfg, Turing.FinTM, Turing.FinTM.bufferTape}
Let $M$ be a finite machine with $k$ work tapes and $x$ an input. Given a
configuration $s$ of $M$ on $x$ (the saved source bank), a captured word $u$, an
optional control state $q$, a position $i \le |x|$, an offset $j \in \mathbb{Z}$,
and an output word $o$, the canonical pairing configuration is the configuration
of the pairing machine over $M$ whose first $k$ work tapes and heads are those of
$s$, whose last tape holds $u$ with head at $j$, whose control is $q$ (halted when
$q$ is absent), whose input head is just before position $i$, and whose output is
$o$.
\end{definition}

\begin{lemma}[Pairing reads the native input]
\lean{Complexity.satPairCfg_input}
\label{Complexity.satPairCfg_input}
\leanok
\uses{Complexity.satPairCfg, Turing.Cfg, Turing.Cfg.inputSymbol, Turing.FinTM, Turing.FinTM.inputSymbol_at}
\difficulty{1}
Let $M$ be a finite machine, $x$ an input, and $i \le |x|$. In every canonical
pairing configuration over $M$ on input $x$ at input position $i$, the symbol read
by the input head is the $i$-th bit of $x$ when $i < |x|$, and the blank when
$i = |x|$.
\end{lemma}

\begin{lemma}[Pairing reads the captured word]
\lean{Complexity.satPairCfg_work}
\label{Complexity.satPairCfg_work}
\leanok
\uses{Complexity.satPairCfg, Turing.Cfg, Turing.Cfg.workTapeSymbols, Turing.FinTM, Turing.FinTM.bufferTape}
\difficulty{1}
Let $M$ be a finite machine, $u$ a captured word, and $j \in \mathbb{Z}$. In every
canonical pairing configuration over $M$ with captured word $u$ and capture head
at offset $j$, the symbol under the last work head is the cell of the buffer tape
of $u$ at $j$, that is, the $j$-th bit of $u$ when $0 \le j < |u|$ and the blank
otherwise.
\end{lemma}

\begin{lemma}[Pairing action on canonical configurations]
\lean{Complexity.satPairAction_apply}
\label{Complexity.satPairAction_apply}
\leanok
\uses{Complexity.satPairAction, Complexity.satPairCfg, Turing.Action.apply, Turing.Cfg, Turing.FinTM, Turing.moveInputPos}
\difficulty{3}
Applying a pairing administrative action to a canonical pairing configuration
changes only the control state, the two displayed head positions, and the output
(appending the optionally emitted bit); the saved source bank and the captured
word are preserved.
\end{lemma}

\begin{lemma}[Rewinding the captured word]
\lean{Complexity.satPair_rewind}
\label{Complexity.satPair_rewind}
\leanok
\uses{Complexity.satPairAction_apply, Complexity.satPairCfg, Complexity.satPairCfg_work, Complexity.satPairTM, Turing.Action.apply, Turing.Cfg, Turing.FinTM, Turing.FinTM.bufferTape, Turing.FinTM.bufferTape_left, Turing.MultiTapeTM.runFrom, Turing.MultiTapeTM.runFrom_succ_eq_step, Turing.MultiTapeTM.runFrom_zero, Turing.MultiTapeTM.step, Turing.moveInputPos_zero}
\difficulty{5}
For every $n \le |u|$: from the canonical pairing configuration in the
capture-rewind state with the capture head at offset $n - 1$ and empty output,
exactly $n + 1$ steps return the head to the origin and enter the replay state,
with all other fields unchanged.
\end{lemma}

\begin{lemma}[Pairing startup seam]
\lean{Complexity.satPair_start}
\label{Complexity.satPair_start}
\leanok
\uses{Complexity.satPairAction, Complexity.satPairAction_apply, Complexity.satPairCfg, Complexity.satPairTM, Complexity.satPair_rewind, Complexity.sat_first_halt, Turing.Action.apply, Turing.Cfg, Turing.Cfg.init, Turing.FinTM, Turing.FinTM.ComputesInTime, Turing.FinTM.bufferTape, Turing.FinTM.timed_rewind, Turing.MultiTapeTM.initCfg, Turing.MultiTapeTM.runFrom, Turing.MultiTapeTM.runFrom_add, Turing.MultiTapeTM.runFrom_succ_eq_step, Turing.MultiTapeTM.step, Turing.captureCfg, Turing.capture_run, Turing.moveInputPos_zero}
\difficulty{6}
Let the source machine $M$ compute the word $u$ on input $x$ within $T$ steps.
Then within $T + |x| + |u| + 5$ steps from its initial configuration, the pairing
machine over $M$ reaches the canonical pairing configuration at the replay entry:
some saved bank of $M$, the captured word $u$ with head at the origin, the input
head at the origin, and all startup output suppressed.
\end{lemma}

\begin{lemma}[Replay with separator]
\lean{Complexity.satPair_replay}
\label{Complexity.satPair_replay}
\leanok
\uses{Complexity.satBits, Complexity.satPairAction_apply, Complexity.satPairCfg, Complexity.satPairCfg_work, Complexity.satPairTM, Turing.Action.apply, Turing.Cfg, Turing.FinTM, Turing.FinTM.bufferTape, Turing.FinTM.bufferTape_nat, Turing.MultiTapeTM.runFrom, Turing.MultiTapeTM.runFrom_succ_eq_step, Turing.MultiTapeTM.step, Turing.moveInputPos_zero}
\difficulty{5}
Let the captured word decompose as $u = p \mathbin{+\!\!+} r$ with the capture
head at offset $|p|$ in the replay state. Then exactly $2\,|r| + 2$ steps emit
the doubling of $r$ followed by the separator $01$, move the capture head to
$|u|$, and enter the native-copy state; the source bank and the native input
head remain fixed.
\end{lemma}

\begin{lemma}[Copying the native suffix]
\lean{Complexity.satPair_copy}
\label{Complexity.satPair_copy}
\leanok
\uses{Complexity.satPairAction_apply, Complexity.satPairCfg, Complexity.satPairCfg_input, Complexity.satPairTM, Turing.Action.apply, Turing.Cfg, Turing.FinTM, Turing.MultiTapeTM.runFrom, Turing.MultiTapeTM.runFrom_succ_eq_step, Turing.MultiTapeTM.runFrom_zero, Turing.MultiTapeTM.step, Turing.moveInputPos_pos_of_ne_right, Turing.moveInputPos_zero}
\difficulty{5}
Let the native input decompose as $x = p \mathbin{+\!\!+} r$ with the pairing
machine in its native-copy state at input position $|p|$. Then exactly $|r| + 1$
steps emit every bit of $r$ and halt on the right blank, leaving all work tapes
and the capture head unchanged.
\end{lemma}

\begin{lemma}[The pairing machine computes the pair]
\lean{Complexity.satPair_computes}
\label{Complexity.satPair_computes}
\leanok
\uses{Complexity.satBits, Complexity.satPairTM, Complexity.satPair_copy, Complexity.satPair_replay, Complexity.satPair_start, Turing.FinTM, Turing.FinTM.ComputesFunInTime, Turing.FinTM.ComputesInTime, Turing.FinTM.ComputesInTime.mono, Turing.FinTM.computesInTime_iff, Turing.MultiTapeTM.output_length_le, Turing.MultiTapeTM.runFrom_add, Turing.pairEncode}
\difficulty{4}
Let the machine $M$ compute a function $f$ within time $T$. Then the pairing
machine over $M$ computes the function $x \mapsto \mathrm{pair}(f(x), x)$ within
time $n \mapsto 4\,T(n) + 2n + 8$: retaining a computed result next to its
original input has a uniform overhead, the capture output length being paid by
the source time.
\end{lemma}

\begin{lemma}[Retaining the input in polynomial time]
\lean{Complexity.sat_pt_pair_input}
\label{Complexity.sat_pt_pair_input}
\leanok
\uses{Complexity.PolyTimeComputable, Complexity.satPairTM, Complexity.satPair_computes, Turing.FinTM.ComputesInTime.mono, Turing.pairEncode}
\difficulty{4}
If a function $f$ on binary strings is polynomial-time computable, then so is the
function $x \mapsto \mathrm{pair}(f(x), x)$ pairing the computed value with the
original input.
\end{lemma}

\begin{lemma}[Mapping the second pair component]
\lean{Complexity.sat_pt_mapSnd}
\label{Complexity.sat_pt_mapSnd}
\leanok
\uses{Complexity.PolyTimeComputable, Turing.FinTM.ComputesInTime.mono, Turing.FinTM.computesFunInTime_pairMapSnd, Turing.pairDecode, Turing.pairEncode}
\difficulty{4}
If a function $g$ on binary strings is polynomial-time computable, then so is the
function that decodes its input as a pair $(a, b)$, returns the pair encoding of
$a$ with $g(b)$, and returns the empty word when the input is not a valid pair
encoding.
\end{lemma}

\begin{lemma}[Pairing two computed fields]
\lean{Complexity.sat_pt_pair}
\label{Complexity.sat_pt_pair}
\leanok
\uses{Complexity.PolyTimeComputable, Complexity.PolyTimeComputable.comp, Complexity.sat_pt_mapSnd, Complexity.sat_pt_pair_input, Turing.pairDecode_pairEncode, Turing.pairEncode}
\difficulty{2}
If $f$ and $g$ are polynomial-time computable functions on binary strings, then
so is $x \mapsto \mathrm{pair}(f(x), g(x))$: independently computed fields can be
paired without losing the original input needed by the second computation.
\end{lemma}

\begin{lemma}[Concatenating two computed words]
\lean{Complexity.sat_pt_append}
\label{Complexity.sat_pt_append}
\leanok
\uses{Complexity.PolyTimeComputable, Complexity.PolyTimeComputable.comp, Complexity.sat_pt_linear, Complexity.sat_pt_pair, Turing.FinTM.computesFunInTime_pairConcat, Turing.pairDecode_pairEncode}
\difficulty{2}
If $f$ and $g$ are polynomial-time computable functions on binary strings, then
so is $x \mapsto f(x) \mathbin{+\!\!+} g(x)$.
\end{lemma}

\begin{definition}[Word-transducer semantics]
\lean{Complexity.satMapWord}
\label{Complexity.satMapWord}
\leanok
Given a state set $S$, a transition $\delta$, a per-bit optional emission
$\varepsilon$, and an optional final emission $\phi$, the transducer semantics is
the word function defined by recursion on the input: on the empty word from state
$q$ it returns the final emission $\phi(q)$ (zero or one bit), and on a bit $b$
followed by a word $r$ from state $q$ it returns the emission $\varepsilon(q, b)$
followed by the semantics of $r$ from state $\delta(q, b)$.
\end{definition}

\begin{definition}[Finite word-transducer machine]
\lean{Complexity.satMapTM}
\label{Complexity.satMapTM}
\leanok
\uses{Turing.FinTM}
Given a finite state set $S$ with a transition $\delta$, a per-bit optional
emission, an optional final emission, and a start state, the word-transducer
machine is the finite Turing machine with no work tapes that, on each input bit,
applies $\delta$, emits the per-bit optional output, and moves right, and at the
right boundary emits the optional final bit and halts.
\end{definition}

\begin{definition}[Transducer configuration]
\lean{Complexity.satMapCfg}
\label{Complexity.satMapCfg}
\leanok
\uses{Turing.Cfg}
For an input word $x$, a control state $q$, a position $i \le |x|$, and an output
word $o$, the transducer configuration is the configuration of a zero-work-tape
machine on input $x$ with control state $q$, input head just before position $i$,
and already-emitted output $o$.
\end{definition}

\begin{lemma}[The transducer machine matches its semantics]
\lean{Complexity.satMap_run}
\label{Complexity.satMap_run}
\leanok
\uses{Complexity.satMapCfg, Complexity.satMapTM, Complexity.satMapWord, Turing.Action.apply, Turing.Cfg.ext_zero_tapes, Turing.Cfg.inputSymbol, Turing.FinTM.inputSymbol_at, Turing.MultiTapeTM.runFrom, Turing.MultiTapeTM.runFrom_succ_eq_step, Turing.MultiTapeTM.step, Turing.moveInputPos_pos_of_ne_right}
\difficulty{5}
Let the input decompose as $x = p \mathbin{+\!\!+} r$. From any state $q$ at
position $|p|$ with emitted prefix $o$, the word-transducer machine halts after
exactly $|r| + 1$ steps with output $o$ followed by the transducer semantics
applied to $q$ and $r$.
\end{lemma}

\begin{lemma}[Transducers are polynomial-time]
\lean{Complexity.satMap_poly}
\label{Complexity.satMap_poly}
\leanok
\uses{Complexity.PolyTimeComputable, Complexity.satMapCfg, Complexity.satMapTM, Complexity.satMapWord, Complexity.satMap_run, Turing.Cfg.ext_zero_tapes, Turing.FinTM.computesInTime_iff, Turing.MultiTapeTM.initCfg}
\difficulty{3}
For every finite state set with transition, per-bit emission, final emission, and
start state, the associated word-transducer function on binary strings is
polynomial-time computable, with the exact time bound $n + 1$.
\end{lemma}

\begin{lemma}[Unary length is polynomial-time]
\lean{Complexity.sat_pt_unaryLength}
\label{Complexity.sat_pt_unaryLength}
\leanok
\uses{Complexity.PolyTimeComputable, Complexity.satMapWord, Complexity.satMap_poly}
\difficulty{3}
The function replacing every bit of its input by $1$ — computing the exact unary
representation of the input length — is polynomial-time computable.
\end{lemma}

\begin{lemma}[Tail is polynomial-time]
\lean{Complexity.sat_pt_tail}
\label{Complexity.sat_pt_tail}
\leanok
\uses{Complexity.PolyTimeComputable, Complexity.satMapWord, Complexity.satMap_poly}
\difficulty{3}
The function removing the first bit of a binary string (returning the empty word
on the empty string) is polynomial-time computable, as a two-state finite
transduction.
\end{lemma}

\begin{lemma}[Head is polynomial-time]
\lean{Complexity.sat_pt_head}
\label{Complexity.sat_pt_head}
\leanok
\uses{Complexity.PolyTimeComputable, Complexity.satMapWord, Complexity.satMap_poly}
\difficulty{3}
The function extracting at most the first bit of a binary string (the length-one
prefix) is polynomial-time computable, as a two-state finite transduction.
\end{lemma}

\begin{lemma}[Equality with a fixed word]
\lean{Complexity.sat_pt_eq}
\label{Complexity.sat_pt_eq}
\leanok
\uses{Complexity.PolyTimeComputable, Complexity.PolyTimeComputable.comp, Complexity.sat_pt_linear, Turing.FinTM.computesFunInTime_ifEq}
\difficulty{3}
Let $f$ be a polynomial-time computable function on binary strings and $w$ a
fixed word. Then the single-bit function testing whether $f(x) = w$ is
polynomial-time computable.
\end{lemma}

\begin{lemma}[Runs agree up to a guarded point]
\lean{Complexity.sat_run_agree}
\label{Complexity.sat_run_agree}
\leanok
\uses{Turing.Cfg, Turing.MultiTapeTM, Turing.MultiTapeTM.runFrom, Turing.MultiTapeTM.runFrom_succ_eq_step', Turing.MultiTapeTM.step}
\difficulty{3}
Let $A$ and $B$ be multi-tape machines with the same tape count and state type,
$c$ a configuration, and $t$ a time such that at every time $i < t$ the one-step
transitions of $B$ and $A$ agree on the configuration reached by $A$ after $i$
steps. Then the $t$-step runs of $B$ and $A$ from $c$ coincide. No eventual halt
is assumed.
\end{lemma}

\begin{lemma}[No early visit to the streaming entry]
\lean{Complexity.satRed_start_guard}
\label{Complexity.satRed_start_guard}
\leanok
\uses{Complexity.satRedTM, Turing.MultiTapeTM.initCfg, Turing.MultiTapeTM.outputSymbol, Turing.MultiTapeTM.output_prefix, Turing.MultiTapeTM.runFrom, Turing.MultiTapeTM.runFrom_succ_eq_step', Turing.MultiTapeTM.step_output}
\difficulty{4}
If the clause-splitting transducer, run for $t$ steps from its initial
configuration, has empty output, then at no earlier time was its control in the
streaming entry state: that state necessarily emits a bit, which output-prefix
monotonicity would make visible at time $t$.
\end{lemma}

\begin{definition}[Intercepted maximum-pass machine]
\lean{Complexity.satMaxTM}
\label{Complexity.satMaxTM}
\leanok
\uses{Complexity.satRedAction, Complexity.satRedTM, Turing.FinTM}
The machine obtained from the clause-splitting transducer by intercepting its
streaming entry state: instead of streaming, it copies the unary counter to the
output (one bit per occupied cell) and halts at the first blank. All other
transitions are taken verbatim from the transducer, so the verified
initialization and maximum pass apply unchanged.
\end{definition}

\begin{lemma}[Interception preserves the startup]
\lean{Complexity.satMax_start}
\label{Complexity.satMax_start}
\leanok
\uses{Complexity.satMaxTM, Complexity.satRedBuffer, Complexity.satRedCfg, Complexity.satRedTM, Complexity.satRed_start, Complexity.satRed_start_guard, Complexity.sat_run_agree, Std.Sat.CNF.numVars, Std.Sat.CNF.serialize, Turing.MultiTapeTM.initCfg, Turing.MultiTapeTM.runFrom, Turing.MultiTapeTM.step}
\difficulty{4}
Let $\varphi$ be a CNF formula with serialization $\mathrm{ser}(\varphi)$ and
its variable bound $N(\varphi)$ (one more than the largest variable index occurring in $\varphi$, or $0$ if no variable occurs). There is a time $t \le 3\,|\mathrm{ser}(\varphi)| + 5$ at which the
intercepted maximum-pass machine, started on $\mathrm{ser}(\varphi)$, is in its
streaming entry state with the counter tape holding $1^{N(\varphi)}$, the buffer
empty apart from its permanent marker, the input head and both work heads at $0$,
and empty output. This is the same endpoint as for the unmodified
clause-splitting transducer.
\end{lemma}

\begin{lemma}[Replaying the counter]
\lean{Complexity.satMax_prefix}
\label{Complexity.satMax_prefix}
\leanok
\uses{Complexity.satMaxTM, Complexity.satRedAction, Complexity.satRedCfg, Complexity.satRedCounter_read, Turing.Action.apply, Turing.Cfg.workTapeSymbols, Turing.MultiTapeTM.runFrom, Turing.MultiTapeTM.runFrom_succ_eq_step', Turing.MultiTapeTM.step, Turing.moveInputPos_zero}
\difficulty{5}
For every $r \le n$: from the streaming entry frame with counter length $n$,
exactly $r$ transitions of the intercepted machine copy the first $r$ unary
counter cells to the output, moving only the counter head; no native input or
buffer movement occurs during the replay.
\end{lemma}

\begin{lemma}[Halting after the counter]
\lean{Complexity.satMax_finish}
\label{Complexity.satMax_finish}
\leanok
\uses{Complexity.satMaxTM, Complexity.satMax_prefix, Complexity.satRedAction, Complexity.satRedCfg, Complexity.satRedCounter_read, Turing.Action.apply, Turing.Cfg.workTapeSymbols, Turing.MultiTapeTM.runFrom, Turing.MultiTapeTM.runFrom_succ_eq_step', Turing.MultiTapeTM.step, Turing.moveInputPos_zero}
\difficulty{4}
Let $x$ be an input, $n \in \mathbb{N}$, and $B$ any buffer tape. Started in the
streaming entry state with the input head at $0$, a unary counter of length $n$
with head at $0$, the buffer $B$ with head at $0$, and empty output, the
intercepted maximum-pass machine after exactly $n + 1$ steps is halted with output
$1^n$ and counter head at $n$, all other fields unchanged; in particular the first
blank after the counter halts the machine without an extra output bit.
\end{lemma}

\begin{lemma}[Computing the unary variable bound]
\lean{Complexity.satMax_computes}
\label{Complexity.satMax_computes}
\leanok
\uses{Complexity.satMaxTM, Complexity.satMax_finish, Complexity.satMax_start, Std.Sat.CNF.decode_serialize, Std.Sat.CNF.numVars, Std.Sat.CNF.numVars_decode_le, Std.Sat.CNF.serialize, Turing.FinTM.ComputesInTime, Turing.FinTM.ComputesInTime.mono, Turing.FinTM.computesInTime_iff, Turing.MultiTapeTM.runFrom_add}
\difficulty{3}
Let $\varphi$ be a CNF formula with serialization $\mathrm{ser}(\varphi)$ and
its variable bound $N(\varphi)$ (one more than the largest variable index occurring in $\varphi$, or $0$ if no variable occurs). The intercepted maximum-pass machine, run on
$\mathrm{ser}(\varphi)$, halts within $4\,|\mathrm{ser}(\varphi)| + 6$ steps, and
its output is the unary word $1^{N(\varphi)}$.
\end{lemma}

\begin{definition}[Input canonicalizer]
\lean{Complexity.satStreamCanonical}
\label{Complexity.satStreamCanonical}
\leanok
\uses{Complexity.satSyntax}
For a binary string $x$, the canonicalization of $x$ is $x$ itself when the
complete syntax scan accepts $x$, and is the one-bit word $0$ (the serialization
of the empty fallback formula) otherwise.
\end{definition}

\begin{lemma}[Canonicalizer specification]
\lean{Complexity.satStreamCanonical_spec}
\label{Complexity.satStreamCanonical_spec}
\leanok
\uses{Complexity.satStreamCanonical, Complexity.satSyntax_spec, Complexity.sat_parse_repr, Std.Sat.CNF.decode, Std.Sat.CNF.fallback, Std.Sat.CNF.parse, Std.Sat.CNF.parse_serialize, Std.Sat.CNF.serialize}
\difficulty{2}
For every binary string $x$, the canonicalization of $x$ equals the serialization
of the formula decoded from $x$ — on every input, including trailing-data
failures.
\end{lemma}

\begin{lemma}[The canonicalizer is polynomial-time]
\lean{Complexity.satStreamCanonical_poly}
\label{Complexity.satStreamCanonical_poly}
\leanok
\uses{Complexity.PolyTimeComputable, Complexity.polyTimeComputable_id, Complexity.satStreamCanonical, Complexity.satSyntax_poly, Complexity.sat_pt_cond, Complexity.sat_pt_const}
\difficulty{1}
The canonicalization function is polynomial-time computable, as a conditional
over the verified complete syntax scanner between the identity and a constant.
\end{lemma}

\begin{lemma}[The unary fresh cursor is polynomial-time]
\lean{Complexity.sat_pt_numVars}
\label{Complexity.sat_pt_numVars}
\leanok
\uses{Complexity.PolyTimeComputable, Complexity.satMaxTM, Complexity.satMax_computes, Complexity.satStreamCanonical, Complexity.satStreamCanonical_poly, Complexity.satStreamCanonical_spec, Complexity.sat_comp_on_image, Std.Sat.CNF.decode, Std.Sat.CNF.numVars, Std.Sat.CNF.serialize, Turing.FinTM.ComputesInTime, Turing.FinTM.ComputesInTime.mono}
\difficulty{5}
For a binary string $x$, let $\varphi_x$ be the CNF formula decoded from $x$ and
$N(\varphi_x)$ its variable bound (one more than the largest variable index
occurring in $\varphi_x$, or $0$ if no variable occurs). The function sending
every binary string $x$ to the unary word $1^{N(\varphi_x)}$ is polynomial-time
computable, on all strings including malformed ones.
\end{lemma}

\begin{lemma}[The encoded startup word is polynomial-time]
\lean{Complexity.satStreamStart_poly}
\label{Complexity.satStreamStart_poly}
\leanok
\uses{Complexity.PolyTimeComputable, Complexity.satStreamStart, Complexity.satStreamWord, Complexity.satSyntax, Complexity.satSyntax_poly, Complexity.sat_pt_cond, Complexity.sat_pt_const, Complexity.sat_pt_numVars, Complexity.sat_pt_pair, Complexity.sat_pt_unaryLength}
\difficulty{3}
The function sending a binary string to the encoded seam word of its startup
state — including the right-boundary offset on invalid inputs — is
polynomial-time computable, before any clause emission.
\end{lemma}
\begin{definition}[Clamped input offset]
\lean{Complexity.satStreamPos}
\label{Complexity.satStreamPos}
\leanok
For an input $x$ and an offset $i \in \mathbb{N}$, the clamped input offset is the
native input-head position at index $\min(i, |x|)$; offsets beyond the right
boundary are clamped to it.
\end{definition}

\begin{lemma}[Reading at a clamped offset]
\lean{Complexity.satStreamPos_read}
\label{Complexity.satStreamPos_read}
\leanok
\uses{Complexity.satStreamPos, Turing.Cfg, Turing.Cfg.inputSymbol, Turing.FinTM.inputSymbol_at}
\difficulty{3}
If a configuration's input position equals the clamped offset $i$ on input $x$,
then its input symbol is the optional $i$-th entry of $x$ (blank beyond the
boundary).
\end{lemma}

\begin{lemma}[Advancing a clamped offset]
\lean{Complexity.satStreamPos_succ}
\label{Complexity.satStreamPos_succ}
\leanok
\uses{Complexity.satStreamPos, Turing.moveInputPos, Turing.moveInputPos_pos_of_ne_right, Turing.moveInputPos_rightBoundary}
\difficulty{4}
Moving the input head right from the clamped offset $i$ yields the clamped offset
$i + 1$ — including repeated requests beyond the native right boundary, which
stay clamped.
\end{lemma}

\begin{definition}[Dropper action]
\lean{Complexity.satDropAction}
\label{Complexity.satDropAction}
\leanok
\uses{Turing.Action}
Given an input-head move $m$, a counter-head move $d$, an optional counter write,
an optional output bit, and an optional next state among six control states, the
dropper action is the one-tape action performing exactly these operations; these
are the actions of the length-controlled suffix extractor.
\end{definition}

\begin{definition}[Suffix-extraction machine]
\lean{Complexity.satDropTM}
\label{Complexity.satDropTM}
\leanok
\uses{Complexity.satDropAction, Turing.FinTM}
The suffix-extraction machine is a six-state finite machine with one work tape.
Given binary words $u$ and $v$ and the input
the pair encoding $\mathrm{pair}(u,v)$ (every bit of $u$ doubled, then the separator $01$, then $v$), it counts the doubled prefix on a unary counter tape,
recognizes the separator, rewinds the counter, skips exactly $|u|$ payload
positions (clamping at the right boundary), and copies the remaining suffix of
$v$. No output precedes recognition of the separator.
\end{definition}

\begin{definition}[Dropper configuration]
\lean{Complexity.satDropCfg}
\label{Complexity.satDropCfg}
\leanok
\uses{Complexity.satRedCounter, Complexity.satStreamPos, Turing.Cfg}
For an input $x$, an optional control state $q$, natural numbers $i$ and $n$, an
offset $h \in \mathbb{Z}$, and an output word $o$, the canonical dropper
configuration is the configuration of the suffix extractor in control $q$, with the
input head at the clamped offset $\min(i, |x|)$, the unary counter tape of length
$n$ with head at $h$, and emitted output $o$.
\end{definition}

\begin{lemma}[Silent dropper action]
\lean{Complexity.satDrop_move}
\label{Complexity.satDrop_move}
\leanok
\uses{Complexity.satDropAction, Complexity.satDropCfg, Complexity.satStreamPos, Turing.Action.apply, Turing.moveInputPos}
\difficulty{2}
Let $x$ be an input, and consider a dropper action with no counter write, input
move $m$, counter-head move $d$, optional output bit $b$, and target state $q'$.
If the move $m$ takes the clamped offset $i$ to the clamped offset $i'$, and
$h + d = h'$, then applying the action to the canonical dropper configuration at
offset $i$ with counter length $n$, counter head $h$, and output $o$ yields the
canonical configuration in state $q'$ at offset $i'$ with the same counter, head
at $h'$, and output $o$ followed by $b$ when present.
\end{lemma}

\begin{lemma}[Extending the counter]
\lean{Complexity.satDrop_write}
\label{Complexity.satDrop_write}
\leanok
\uses{Complexity.satDropAction, Complexity.satDropCfg, Complexity.satRedCounter_write, Complexity.satStreamPos_succ, Turing.Action.apply}
\difficulty{3}
Let $x$ be an input, $i$ a natural offset, and $n \in \mathbb{N}$. Applying the
dropper action that writes a one under the counter head, moves the input head and
the counter head right, and enters the counting state, to a canonical dropper
configuration on $x$ at clamped offset $i$ whose unary counter has length $n$ and
head at $n$ (with empty output), yields the canonical configuration in the
counting state at clamped offset $i + 1$ with counter length $n + 1$ and counter
head at $n + 1$: writing the counter's right blank appends exactly one unary cell.
\end{lemma}

\begin{lemma}[Counting one doubled bit]
\lean{Complexity.satDrop_double}
\label{Complexity.satDrop_double}
\leanok
\uses{Complexity.satDropCfg, Complexity.satDropTM, Complexity.satDrop_move, Complexity.satDrop_write, Complexity.satStreamPos_read, Complexity.satStreamPos_succ, Turing.Action.apply, Turing.Cfg.inputSymbol, Turing.MultiTapeTM.runFrom, Turing.MultiTapeTM.step}
\difficulty{4}
Let $x$ be an input of the form $p$, followed by two copies of a bit $b$, followed
by a word $r$, and let $n \in \mathbb{N}$. From the counting state at offset $|p|$
with a unary counter of length $n$, counter head at $n$, and empty output, two
steps of the suffix extractor yield the counting state at offset $|p| + 2$ with
counter length $n + 1$ and counter head at $n + 1$; nothing is emitted.
\end{lemma}

\begin{lemma}[Parsing the doubled prefix]
\lean{Complexity.satDrop_parse}
\label{Complexity.satDrop_parse}
\leanok
\uses{Complexity.satBits, Complexity.satBits_length, Complexity.satDropCfg, Complexity.satDropTM, Complexity.satDrop_double, Turing.MultiTapeTM.runFrom, Turing.MultiTapeTM.runFrom_add, Turing.MultiTapeTM.runFrom_zero, Turing.pairEncode}
\difficulty{4}
Let $u$ and $v$ be binary words, with $u = p \mathbin{+\!\!+} r$ split into a
prefix $p$ and a suffix $r$, and consider the input
the pair encoding $\mathrm{pair}(u,v)$ (every bit of $u$ doubled, then the separator $01$, then $v$). Started in its counting state at input offset $2|p|$, with a
unary counter of length $|p|$ and counter head at $|p|$ and empty output, the
suffix extractor after exactly $2|r|$ steps is again in its counting state, at
input offset $2|u|$, with a unary counter of length $|u|$ and counter head at
$|u|$; nothing is emitted.
\end{lemma}

\begin{lemma}[Recognizing the separator]
\lean{Complexity.satDrop_separator}
\label{Complexity.satDrop_separator}
\leanok
\uses{Complexity.satBits, Complexity.satBits_length, Complexity.satDropCfg, Complexity.satDropTM, Complexity.satDrop_move, Complexity.satStreamPos_read, Complexity.satStreamPos_succ, Turing.Action.apply, Turing.Cfg.inputSymbol, Turing.MultiTapeTM.runFrom, Turing.MultiTapeTM.step, Turing.pairEncode}
\difficulty{4}
Let $u$ and $v$ be binary words and consider the input
the pair encoding $\mathrm{pair}(u,v)$ (every bit of $u$ doubled, then the separator $01$, then $v$). Started in its counting state at input offset $2|u|$, with a
unary counter of length $|u|$, counter head at $|u|$, and empty output, the suffix
extractor after exactly two steps is in its counter-rewind state at input offset
$2|u| + 2$ (the first payload bit, just past the separator), with the counter
unchanged, counter head at $|u| - 1$, and nothing emitted.
\end{lemma}

\begin{lemma}[Dropper counter rewind]
\lean{Complexity.satDrop_rewind}
\label{Complexity.satDrop_rewind}
\leanok
\uses{Complexity.satDropCfg, Complexity.satDropTM, Complexity.satDrop_move, Complexity.satRedCounter_left, Complexity.satRedCounter_read, Turing.Action.apply, Turing.Cfg.workTapeSymbols, Turing.MultiTapeTM.runFrom, Turing.MultiTapeTM.runFrom_succ_eq_step, Turing.MultiTapeTM.runFrom_zero, Turing.MultiTapeTM.step, Turing.moveInputPos_zero}
\difficulty{4}
For every $r \le n$: from the rewind state with counter length $n$ and counter
head at offset $r - 1$, exactly $r + 1$ steps return the counter head to $0$ and
enter the skip state, without disturbing the payload head or emitting.
\end{lemma}

\begin{lemma}[Skipping by the stored length]
\lean{Complexity.satDrop_skip}
\label{Complexity.satDrop_skip}
\leanok
\uses{Complexity.satDropCfg, Complexity.satDropTM, Complexity.satDrop_move, Complexity.satRedCounter_read, Complexity.satStreamPos_succ, Turing.Action.apply, Turing.Cfg.workTapeSymbols, Turing.MultiTapeTM.runFrom, Turing.MultiTapeTM.runFrom_succ_eq_step', Turing.MultiTapeTM.runFrom_zero, Turing.MultiTapeTM.step}
\difficulty{4}
For every $r \le n$: in the skip state with counter length $n$, $r$ steps advance
the clamped payload offset by $r$ while walking the counter head to $r$. Native
clamping makes this total even when the stored length exceeds the payload.
\end{lemma}

\begin{lemma}[Dispatch to the copy phase]
\lean{Complexity.satDrop_dispatch}
\label{Complexity.satDrop_dispatch}
\leanok
\uses{Complexity.satDropCfg, Complexity.satDropTM, Complexity.satDrop_move, Complexity.satRedCounter_read, Turing.Action.apply, Turing.Cfg.workTapeSymbols, Turing.MultiTapeTM.step, Turing.moveInputPos_zero}
\difficulty{3}
Let $x$ be an input and $p, n \in \mathbb{N}$. From the skip state at clamped
offset $p$, with a unary counter of length $n$ whose head is at the first blank
$n$ and empty output, one step of the suffix extractor yields the configuration in
the copy state with every other field unchanged; in particular the input head does
not move past the first unconsumed payload bit.
\end{lemma}

\begin{lemma}[Copying the remaining suffix]
\lean{Complexity.satDrop_copy}
\label{Complexity.satDrop_copy}
\leanok
\uses{Complexity.satDropCfg, Complexity.satDropTM, Complexity.satDrop_move, Complexity.satStreamPos_read, Complexity.satStreamPos_succ, Turing.Action.apply, Turing.Cfg.inputSymbol, Turing.MultiTapeTM.runFrom, Turing.MultiTapeTM.runFrom_succ_eq_step, Turing.MultiTapeTM.runFrom_zero, Turing.MultiTapeTM.step, Turing.moveInputPos_zero}
\difficulty{5}
Let the native input decompose as $x = p \mathbin{+\!\!+} r$ with the suffix
extractor in its copy state at clamped offset $|p|$. Then exactly $|r| + 1$ steps
emit every bit of $r$ and halt on the right blank, with the counter untouched.
\end{lemma}

\begin{lemma}[Linear-time suffix extraction]
\lean{Complexity.satDrop_computes}
\label{Complexity.satDrop_computes}
\leanok
\uses{Complexity.satBits, Complexity.satBits_length, Complexity.satDropCfg, Complexity.satDropTM, Complexity.satDrop_copy, Complexity.satDrop_dispatch, Complexity.satDrop_parse, Complexity.satDrop_rewind, Complexity.satDrop_separator, Complexity.satDrop_skip, Complexity.satRedCounter, Complexity.satStreamPos, Turing.Cfg.init, Turing.FinTM.ComputesInTime, Turing.FinTM.ComputesInTime.mono, Turing.FinTM.bufferTape, Turing.FinTM.computesInTime_iff, Turing.MultiTapeTM.initCfg, Turing.MultiTapeTM.runFrom, Turing.MultiTapeTM.runFrom_add, Turing.MultiTapeTM.runFrom_succ_eq_step', Turing.pairEncode}
\difficulty{6}
For all binary words $u$ and $v$, the suffix extractor run on
the pair encoding $\mathrm{pair}(u,v)$ (every bit of $u$ doubled, then the separator $01$, then $v$) halts within $3(\ell + 1)$ steps, where $\ell$ is the length of
that paired input, and its output is the suffix of $v$ obtained by dropping its
first $|u|$ bits.
\end{lemma}

\begin{lemma}[Pair projections are polynomial-time]
\lean{Complexity.sat_pt_fields}
\label{Complexity.sat_pt_fields}
\leanok
\uses{Complexity.PolyTimeComputable, Complexity.satStreamFst, Complexity.satStreamSnd, Complexity.sat_pt_linear, Turing.FinTM.computesFunInTime_pairFst, Turing.FinTM.computesFunInTime_pairSnd}
\difficulty{1}
Both total pair projections — first and second component of a pair-encoded word,
defaulting to the empty word on malformed inputs — are polynomial-time
computable.
\end{lemma}

\begin{lemma}[Projected suffix extraction is polynomial-time]
\lean{Complexity.sat_pt_drop}
\label{Complexity.sat_pt_drop}
\leanok
\uses{Complexity.PolyTimeComputable, Complexity.satDropTM, Complexity.satDrop_computes, Complexity.satStreamFst, Complexity.satStreamSnd, Complexity.sat_comp_on_image, Complexity.sat_pt_fields, Complexity.sat_pt_pair, Turing.FinTM.ComputesInTime, Turing.FinTM.ComputesInTime.mono, Turing.FinTM.computesInTime_iff, Turing.MultiTapeTM.output_length_le, Turing.pairEncode}
\difficulty{5}
The function sending a word $z$ to the suffix of its second pair projection
obtained by dropping as many bits as the length of its first pair projection is
polynomial-time computable on every word, malformed pair inputs being interpreted
through the empty-default projections.
\end{lemma}

\begin{lemma}[Token split of a serialized literal]
\lean{Complexity.satToken_literal}
\label{Complexity.satToken_literal}
\leanok
\uses{Std.Sat.CNF.serializeLit, Turing.unaryTokenSplit}
\difficulty{3}
For every literal $\ell = (v, b)$ and remainder $r$, the unary token splitter
applied to the serialization of $\ell$ followed by $r$ returns the token
$1^{v+1}0$ together with the suffix $b :: r$: the delimiter is consumed and the
polarity bit remains the first bit of the returned suffix.
\end{lemma}

\begin{lemma}[Shape of the token split]
\lean{Complexity.satToken_shape}
\label{Complexity.satToken_shape}
\leanok
\uses{Std.Sat.CNF.takeTrues, Turing.unaryTokenSplit}
\difficulty{3}
For every binary string $x$: the unary token splitter returns the leading run of
ones together with at most one following bit as the token, and the rest as the
suffix; moreover the residue after the leading run never starts with another
one.
\end{lemma}

\begin{lemma}[Token split on a failed literal parse]
\lean{Complexity.satToken_failure}
\label{Complexity.satToken_failure}
\leanok
\uses{Complexity.satToken_shape, Std.Sat.CNF.parseLit, Std.Sat.CNF.takeTrues, Turing.unaryTokenSplit}
\difficulty{4}
If the literal parser fails on a string starting with a one (a literal marker),
then the unary token splitter on that string returns an empty suffix: there is no
polarity bit, which keeps the total round implementation faithful even on locally
malformed state--input combinations.
\end{lemma}

\begin{definition}[Round request: fresh field]
\lean{Complexity.satReqFresh}
\label{Complexity.satReqFresh}
\leanok
\uses{Complexity.satStreamFst}
A round request is the pair encoding of an encoded seam word with a fresh copy of
the native input. For a word $z$, the fresh field of $z$ is the first component of
the first component of $z$ (each projection defaulting to the empty word on
malformed inputs); on a canonical request it is the unary fresh cursor.
\end{definition}

\begin{definition}[Round request: consumed field]
\lean{Complexity.satReqUsed}
\label{Complexity.satReqUsed}
\leanok
\uses{Complexity.satStreamFst, Complexity.satStreamSnd}
For a word $z$, the consumed field of $z$ is the first component of the second
component of the first component of $z$ (each projection defaulting to the empty
word on malformed inputs); on a canonical round request it is the unary consumed
length.
\end{definition}

\begin{definition}[Round request: phase field]
\lean{Complexity.satReqPhase}
\label{Complexity.satReqPhase}
\leanok
\uses{Complexity.satStreamFst, Complexity.satStreamSnd}
For a word $z$, the phase field of $z$ is the second component of the second
component of the first component of $z$ (each projection defaulting to the empty
word on malformed inputs); on a canonical round request it is the unary phase tag.
\end{definition}

\begin{definition}[Round request: unread input]
\lean{Complexity.satReqRest}
\label{Complexity.satReqRest}
\leanok
\uses{Complexity.satReqUsed, Complexity.satStreamSnd}
For a word $z$, the unread input of $z$ is the second component of $z$ with as
many leading bits dropped as the length of the consumed field of $z$.
\end{definition}

\begin{definition}[Round request: polarity suffix]
\lean{Complexity.satReqPol}
\label{Complexity.satReqPol}
\leanok
\uses{Complexity.satReqRest, Turing.unaryTokenSplit}
The suffix returned by the unary token splitter on the unread input of a round
request; its first bit, when present, is the polarity bit of the literal being
read.
\end{definition}

\begin{definition}[Round request: current literal]
\lean{Complexity.satReqLit}
\label{Complexity.satReqLit}
\leanok
\uses{Complexity.satReqPol, Complexity.satReqRest, Turing.unaryTokenSplit}
For a word $z$, the current literal of $z$ is the token returned by the unary
token splitter on the unread input of $z$, extended by at most the first bit of
the polarity suffix of $z$.
\end{definition}

\begin{definition}[Tail-link decision]
\lean{Complexity.satReqLink}
\label{Complexity.satReqLink}
\leanok
\uses{Complexity.satReqPhase, Complexity.satReqPol}
The Boolean deciding whether the current round performs a fresh-variable tail
link: the phase field equals the tail tag ($1^3$) and the bit after the polarity
bit is a one (another literal follows in the clause).
\end{definition}

\begin{definition}[Packing the canonical fields]
\lean{Complexity.satReqPack}
\label{Complexity.satReqPack}
\leanok
\uses{Turing.pairEncode}
The canonical packing of a fresh field, a consumed field, and a fixed phase tag
$q$ into an encoded seam word: the first field paired with the pair of the second
field and the unary word $1^q$. No round-local scratch is involved.
\end{definition}

\begin{definition}[Tail-link fragment]
\lean{Complexity.satReqFragment}
\label{Complexity.satReqFragment}
\leanok
\uses{Complexity.satReqFresh, Complexity.satReqLit}
The one larger output chunk of a tail-link round, assembled from the request
fields: the positive fresh link ($1$ followed by the unary fresh field, then
$01$), two clause markers, the negative fresh link, and the consumed original
literal, in that order.
\end{definition}

\begin{definition}[Word-level chunk program]
\lean{Complexity.satReqEmit}
\label{Complexity.satReqEmit}
\leanok
\uses{Complexity.satReqFragment, Complexity.satReqLink, Complexity.satReqLit, Complexity.satReqPhase, Complexity.satReqPol, Complexity.satReqRest}
A total straight-line word program computing the chunk emitted by one streaming
round from a round request: finished phase emits nothing; exhausted input or a
terminator bit emits $0$; a clause marker at formula phase emits $1$; a missing
polarity bit emits nothing; a tail link emits the assembled fragment; otherwise
the current literal is copied. Only the grammar phase and the complete-literal
lookahead are tested locally — full validation belongs to startup.
\end{definition}

\begin{definition}[Word-level next-state program]
\lean{Complexity.satReqStep}
\label{Complexity.satReqStep}
\leanok
\uses{Complexity.satReqFresh, Complexity.satReqLink, Complexity.satReqLit, Complexity.satReqPack, Complexity.satReqPhase, Complexity.satReqPol, Complexity.satReqRest, Complexity.satReqUsed, Complexity.satStreamFst}
A total straight-line word program installing the next encoded seam word from a
round request, by the same local tests as the chunk program: the consumed field
advances by the whole current literal length (or by one for markers and
terminators), the phase tag is updated, and the fresh field grows by one exactly
on a tail link.
\end{definition}

\begin{lemma}[First-bit equality test]
\lean{Complexity.sat_take_one_true}
\label{Complexity.sat_take_one_true}
\leanok
\difficulty{1}
For every binary string $x$, the length-one prefix of $x$ equals $[1]$ if and
only if the optional first bit of $x$ is a one.
\end{lemma}

\begin{lemma}[Recovering the request fields]
\lean{Complexity.satReq_fields}
\label{Complexity.satReq_fields}
\leanok
\uses{Complexity.SatStreamState, Complexity.satReqFresh, Complexity.satReqPhase, Complexity.satReqRest, Complexity.satReqUsed, Complexity.satStreamFst, Complexity.satStreamSnd, Complexity.satStreamWord, Turing.pairDecode_pairEncode, Turing.pairEncode}
\difficulty{1}
Let $x$ be a binary word and $s$ a seam state with fresh cursor $f$, consumed
length $c$, and phase $q$, and let $w_s$ be the encoded seam word of $s$. On the
canonical round request $\mathrm{pair}(w_s, x)$, the fresh, consumed, phase, and
unread-input projections return exactly $1^f$, $1^c$, $1^q$, and the word $x$ with
its first $c$ bits dropped.
\end{lemma}

\begin{lemma}[The word programs implement the round]
\lean{Complexity.satReq_round}
\label{Complexity.satReq_round}
\leanok
\uses{Complexity.SatStreamState, Complexity.satReqEmit, Complexity.satReqFragment, Complexity.satReqFresh, Complexity.satReqLink, Complexity.satReqLit, Complexity.satReqPack, Complexity.satReqPhase, Complexity.satReqPol, Complexity.satReqRest, Complexity.satReqStep, Complexity.satReqUsed, Complexity.satReq_fields, Complexity.satStreamFst, Complexity.satStreamRound, Complexity.satStreamSnd, Complexity.satStreamWord, Complexity.satToken_failure, Complexity.satToken_literal, Complexity.sat_parseLit_repr, Complexity.sat_take_one_true, Std.Sat.CNF.parseLit, Std.Sat.CNF.serializeLit, Turing.pairDecode_pairEncode, Turing.pairEncode}
\difficulty{6}
On every canonical round request — the pair encoding of the encoded seam word of
any state $s$ (not only reachable ones) with the input $x$ — the word-level
next-state program returns the encoded word of the state after one streaming
round, and the word-level chunk program returns exactly the chunk emitted by that
round.
\end{lemma}

\begin{lemma}[Request fields are polynomial-time]
\lean{Complexity.satReq_fields_poly}
\label{Complexity.satReq_fields_poly}
\leanok
\uses{Complexity.PolyTimeComputable, Complexity.PolyTimeComputable.comp, Complexity.satReqFresh, Complexity.satReqLit, Complexity.satReqPhase, Complexity.satReqPol, Complexity.satReqRest, Complexity.satReqUsed, Complexity.satStreamFst, Complexity.satStreamSnd, Complexity.sat_pt_append, Complexity.sat_pt_drop, Complexity.sat_pt_fields, Complexity.sat_pt_head, Complexity.sat_pt_linear, Complexity.sat_pt_pair, Turing.FinTM.computesFunInTime_unaryToken, Turing.pairDecode_pairEncode, Turing.pairEncode, Turing.unaryTokenSplit}
\difficulty{4}
All six read-only request projections — fresh field, consumed field, phase
field, unread input, polarity suffix, and current literal — are polynomial-time
computable.
\end{lemma}

\begin{lemma}[Packing is polynomial-time]
\lean{Complexity.sat_pt_pack}
\label{Complexity.sat_pt_pack}
\leanok
\uses{Complexity.PolyTimeComputable, Complexity.satReqPack, Complexity.sat_pt_const, Complexity.sat_pt_pair}
\difficulty{1}
If two word functions are polynomial-time computable, then for every fixed phase
tag so is the function packing their values into an encoded seam word.
\end{lemma}

\begin{lemma}[The link decision is polynomial-time]
\lean{Complexity.satReqLink_poly}
\label{Complexity.satReqLink_poly}
\leanok
\uses{Complexity.PolyTimeComputable, Complexity.PolyTimeComputable.comp, Complexity.satReqLink, Complexity.satReq_fields_poly, Complexity.sat_pt_and, Complexity.sat_pt_eq, Complexity.sat_pt_head, Complexity.sat_pt_tail}
\difficulty{2}
The single-bit tail-link decision on round requests is polynomial-time
computable, as a conjunction of two finite-word equality tests.
\end{lemma}

\begin{lemma}[The fragment is polynomial-time]
\lean{Complexity.satReqFragment_poly}
\label{Complexity.satReqFragment_poly}
\leanok
\uses{Complexity.PolyTimeComputable, Complexity.PolyTimeComputable.comp, Complexity.satReqFragment, Complexity.satReqFresh, Complexity.satReq_fields_poly, Complexity.sat_pt_append, Complexity.sat_pt_const, Complexity.sat_pt_linear, Turing.FinTM.computesFunInTime_prepend}
\difficulty{3}
The tail-link fragment function on round requests is polynomial-time computable,
built by a fixed number of polynomial-time concatenations.
\end{lemma}

\begin{lemma}[The chunk program is polynomial-time]
\lean{Complexity.satReqEmit_poly}
\label{Complexity.satReqEmit_poly}
\leanok
\uses{Complexity.PolyTimeComputable, Complexity.PolyTimeComputable.comp, Complexity.satReqEmit, Complexity.satReqFragment_poly, Complexity.satReqLink_poly, Complexity.satReq_fields_poly, Complexity.sat_pt_cond, Complexity.sat_pt_const, Complexity.sat_pt_eq, Complexity.sat_pt_head}
\difficulty{4}
The word-level per-round chunk program is polynomial-time computable — including
the empty finished chunks and the invalid-input fallback — with an actual finite
machine witness.
\end{lemma}

\begin{lemma}[The next-state program is polynomial-time]
\lean{Complexity.satReqStep_poly}
\label{Complexity.satReqStep_poly}
\leanok
\uses{Complexity.PolyTimeComputable, Complexity.PolyTimeComputable.comp, Complexity.satReqLink_poly, Complexity.satReqStep, Complexity.satReq_fields_poly, Complexity.sat_pt_append, Complexity.sat_pt_cond, Complexity.sat_pt_const, Complexity.sat_pt_eq, Complexity.sat_pt_fields, Complexity.sat_pt_head, Complexity.sat_pt_pack, Complexity.sat_pt_unaryLength}
\difficulty{4}
The word-level next-state program is polynomial-time computable; all branch
tests are computed and captured on the original round request.
\end{lemma}
\begin{definition}[Marker-free appender machine]
\lean{Complexity.satAppendTM}
\label{Complexity.satAppendTM}
\leanok
\uses{Turing.FinTM, Turing.FinTM.controlAction}
A five-state, one-tape administrative machine that appends the native input to
the word already on its work tape and then rewinds both heads: seek the first
blank after the existing word, copy the input, rewind the work head over the
occupied prefix, and rewind the input head. The only occupied tape is the data
tape; no origin marker is introduced. It is used to form a round request and at
startup of the streaming host.
\end{definition}

\begin{definition}[Appender configuration]
\lean{Complexity.satAppendCfg}
\label{Complexity.satAppendCfg}
\leanok
\uses{Turing.Cfg, Turing.FinTM.bufferTape}
The complete configuration of the appender: a control state, an input position,
the buffer tape of a stated word with its head at a stated offset, and empty
output; there is no hidden scratch.
\end{definition}

\begin{lemma}[Seeking the first blank]
\lean{Complexity.satAppend_seek}
\label{Complexity.satAppend_seek}
\leanok
\uses{Complexity.satAppendCfg, Complexity.satAppendTM, Turing.Action.apply, Turing.Cfg.workTapeSymbols, Turing.MultiTapeTM.runFrom, Turing.MultiTapeTM.runFrom_succ_eq_step, Turing.MultiTapeTM.runFrom_zero, Turing.MultiTapeTM.step}
\difficulty{4}
Let the work tape hold a word $w = p \mathbin{+\!\!+} r$ with the head at offset
$|p|$ in the seek state. Then exactly $|r| + 1$ steps bring the head to the first
blank at offset $|w|$ and enter the copy state, without moving the native input
head.
\end{lemma}

\begin{lemma}[Appending the native input]
\lean{Complexity.satAppend_copy}
\label{Complexity.satAppend_copy}
\leanok
\uses{Complexity.satAppendCfg, Complexity.satAppendTM, Complexity.satStreamPos, Complexity.satStreamPos_read, Complexity.satStreamPos_succ, Turing.Action.apply, Turing.Cfg.inputSymbol, Turing.FinTM.bufferTape_append, Turing.MultiTapeTM.runFrom, Turing.MultiTapeTM.runFrom_succ_eq_step, Turing.MultiTapeTM.runFrom_zero, Turing.MultiTapeTM.step, Turing.moveInputPos_zero}
\difficulty{5}
Let the native input decompose as $x = p \mathbin{+\!\!+} r$, with the appender
in its copy state at clamped input offset $|p|$ and work tape holding
$w \mathbin{+\!\!+} p$ with the head at its right blank. Then exactly $|r| + 1$
steps append the remaining suffix $r$ to the work tape and enter the work-tape
rewind with the head on the last occupied cell of $w \mathbin{+\!\!+} x$.
\end{lemma}

\begin{lemma}[Appender work rewind]
\lean{Complexity.satAppend_rewind}
\label{Complexity.satAppend_rewind}
\leanok
\uses{Complexity.satAppendCfg, Complexity.satAppendTM, Turing.Action.apply, Turing.Cfg.workTapeSymbols, Turing.FinTM.bufferTape, Turing.FinTM.bufferTape_nat, Turing.MultiTapeTM.runFrom, Turing.MultiTapeTM.runFrom_succ_eq_step, Turing.MultiTapeTM.runFrom_zero, Turing.MultiTapeTM.step}
\difficulty{4}
For every $r \le |w|$: from the rewind state with the work head at offset
$r - 1$ over a tape holding $w$, exactly $r + 1$ steps return the head to the
origin and enter the input-rewind state; the blank at offset $-1$ is used only
locally and never written.
\end{lemma}

\begin{lemma}[Cutting a padded run at its first halt]
\lean{Complexity.sat_call_first_halt}
\label{Complexity.sat_call_first_halt}
\leanok
\uses{Turing.Cfg, Turing.Cfg.Halted, Turing.MultiTapeTM, Turing.MultiTapeTM.runFrom, Turing.MultiTapeTM.runFrom_add, Turing.MultiTapeTM.runFrom_of_halt, Turing.MultiTapeTM.runFrom_zero}
\difficulty{4}
Let a multi-tape machine, started in a live configuration, be halted after $t$
steps. Then there is a positive time $u \le t$ such that the machine is not
halted at any earlier time, and the $u$-step run equals the $t$-step run in
every field: a padded halting run can be cut at its first halt without changing
any endpoint.
\end{lemma}

\begin{lemma}[Clean return of the appender]
\lean{Complexity.satAppend_clean}
\label{Complexity.satAppend_clean}
\leanok
\uses{Complexity.satAppendCfg, Complexity.satAppendTM, Complexity.satAppend_copy, Complexity.satAppend_rewind, Complexity.satAppend_seek, Complexity.satStreamPos, Complexity.sat_call_first_halt, Turing.Cfg.Halted, Turing.Cfg.ofWords, Turing.FinTM.timed_rewind, Turing.MultiTapeTM.runFrom, Turing.MultiTapeTM.runFrom_add, Turing.stateWord}
\difficulty{5}
For every native input $x$ and existing word $w$ on the data tape, there is a
positive time $t \le 2|w| + 3|x| + 6$ such that the appender, started at its
canonical entry over the state word of $w$, is not halted before $t$ and at time
$t$ halts in the canonical configuration holding the state word of
$w \mathbin{+\!\!+} x$, with both the input and work heads returned to their
origins and no output.
\end{lemma}

\begin{definition}[Padded action]
\lean{Complexity.satPadAction}
\label{Complexity.satPadAction}
\leanok
\uses{Turing.Action}
The embedding of an action over $k$ work tapes into an action over $K \ge k$
work tapes: the extra tapes are neither written nor moved.
\end{definition}

\begin{definition}[Padded machine]
\lean{Complexity.satPadTM}
\label{Complexity.satPadTM}
\leanok
\uses{Complexity.satPadAction, Turing.FinTM}
The embedding of a finite machine with $k$ work tapes into one with $K \ge k$
work tapes by padding every transition; all modules of the streaming host share
tape zero, with each source bank padded on the right.
\end{definition}

\begin{definition}[Padded configuration]
\lean{Complexity.satPadCfg}
\label{Complexity.satPadCfg}
\leanok
\uses{Turing.Cfg}
The embedding of a configuration over $k$ work tapes into one over $K \ge k$
work tapes: the padding tapes are genuinely blank with heads at the origin.
\end{definition}

\begin{lemma}[Padding commutes with actions]
\lean{Complexity.satPad_apply}
\label{Complexity.satPad_apply}
\leanok
\uses{Complexity.satPadAction, Complexity.satPadCfg, Turing.Action, Turing.Action.apply, Turing.Cfg}
\difficulty{3}
Applying the padded version of an action to the padded version of a
configuration equals padding the result of applying the original action —
including writes at arbitrary cells.
\end{lemma}

\begin{lemma}[Padding commutes with runs]
\lean{Complexity.satPad_run}
\label{Complexity.satPad_run}
\leanok
\uses{Complexity.satPadAction, Complexity.satPadCfg, Complexity.satPadTM, Complexity.satPad_apply, Turing.Action.apply, Turing.Cfg, Turing.Cfg.inputSymbol, Turing.Cfg.workTapeSymbols, Turing.FinTM, Turing.MultiTapeTM.runFrom, Turing.MultiTapeTM.runFrom_comm_of_step, Turing.MultiTapeTM.step}
\difficulty{4}
For every machine $M$, tape bound $K \ge k$, configuration $c$, and time $t$:
running the padded machine for $t$ steps from the padded configuration equals
padding the $t$-step run of $M$ from $c$ — the padded source observes exactly
its original work symbols.
\end{lemma}

\begin{lemma}[Padded canonical seams]
\lean{Complexity.satPad_seam}
\label{Complexity.satPad_seam}
\leanok
\uses{Complexity.satPadCfg, Turing.Cfg.ofWords, Turing.FinTM.bufferTape, Turing.stateWord}
\difficulty{3}
For a positive source tape count, padding the canonical configuration built from
a state word (data word on tape zero, other tapes blank) yields the canonical
configuration over the larger tape count built from the same state word, with
the same output.
\end{lemma}

\begin{definition}[Call-stopping wrapper]
\lean{Complexity.satStopTM}
\label{Complexity.satStopTM}
\leanok
\uses{Turing.FinTM, Turing.FinTM.controlAction}
Given a machine $C$ with designated entry and exit states, the wrapper machine
whose states carry an additional Boolean release flag: it simulates $C$, raising
the flag after the first step, and halts silently as soon as the flag is raised
and the control is at the designated exit. The flag ensures that even an entry
equal to the exit executes at least once.
\end{definition}

\begin{definition}[Release-flag embedding]
\lean{Complexity.satStopCfg}
\label{Complexity.satStopCfg}
\leanok
\uses{Turing.Cfg}
The embedding of a configuration of a machine into the call-stopping wrapper
with a given release-flag value; all physical fields (tapes, heads, output) are
preserved.
\end{definition}

\begin{lemma}[Stopped step before the exit]
\lean{Complexity.satStop_step}
\label{Complexity.satStop_step}
\leanok
\uses{Complexity.satStopCfg, Complexity.satStopTM, Turing.Cfg, Turing.FinTM, Turing.MultiTapeTM.step}
\difficulty{3}
If the embedded configuration is live and not (with the flag raised) at the
designated exit, then one step of the call-stopping wrapper performs the same
physical step as the wrapped machine and raises the release flag.
\end{lemma}

\begin{lemma}[Liveness of run prefixes]
\lean{Complexity.sat_live_prefix}
\label{Complexity.sat_live_prefix}
\leanok
\uses{Turing.Cfg, Turing.MultiTapeTM, Turing.MultiTapeTM.runFrom, Turing.MultiTapeTM.runFrom_add, Turing.MultiTapeTM.runFrom_of_halt}
\difficulty{3}
If a multi-tape machine is live (not halted) after $t$ steps from a
configuration, then it is live after every earlier number of steps.
\end{lemma}

\begin{lemma}[Stopping a clean call]
\lean{Complexity.satStop_clean}
\label{Complexity.satStop_clean}
\leanok
\uses{Complexity.satStopCfg, Complexity.satStopTM, Complexity.satStop_step, Complexity.sat_live_prefix, Turing.Action.apply, Turing.Cfg.ofWords, Turing.FinTM, Turing.FinTM.controlAction, Turing.MultiTapeTM.runFrom, Turing.MultiTapeTM.runFrom_succ_eq_step', Turing.MultiTapeTM.step, Turing.stateWord}
\difficulty{5}
Let a machine $C$, started at its entry over the state word of an argument,
avoid its designated exit at all intermediate positive times before $t > 0$ and
reach at time $t$ the exit over the state word of a result with a given output.
Then the call-stopping wrapper, started with the flag down, runs for $t + 1$
steps and halts in that same physical configuration (flag raised, state halted),
retaining every emitted bit.
\end{lemma}

\begin{definition}[Host state space]
\lean{Complexity.SatHostState}
\label{Complexity.SatHostState}
\leanok
\uses{Turing.FinTM}
The control-state space of the streaming host: a single anchor state plus the
disjoint union of the state spaces of six finite modules (startup append and
install, then the recurring pack, append, emit, and install calls).
\end{definition}

\begin{definition}[Module return successor]
\lean{Complexity.satHostNext}
\label{Complexity.satHostNext}
\leanok
The return destination of each of the six modules: the startup appender
continues to the startup installer, and the pack, append, and emit calls
continue to the next module; the two installers return to the unique loop
anchor.
\end{definition}

\begin{definition}[Module entry state]
\lean{Complexity.satHostEntry}
\label{Complexity.satHostEntry}
\leanok
\uses{Complexity.SatHostState, Turing.FinTM}
Given six finite modules and an index $i \in \{0, \dots, 5\}$, the entry state of
module $i$ is the host control state corresponding to the initial state of the
$i$-th module inside the finite body.
\end{definition}

\begin{definition}[Module return state]
\lean{Complexity.satHostRet}
\label{Complexity.satHostRet}
\leanok
\uses{Complexity.SatHostState, Complexity.satHostEntry, Complexity.satHostNext, Turing.FinTM}
The host control state to which the $i$-th module returns: the entry of its
successor module, or the loop anchor when there is none. No additional tape
operation is performed on return.
\end{definition}

\begin{definition}[The streaming host machine]
\lean{Complexity.satHostTM}
\label{Complexity.satHostTM}
\leanok
\uses{Complexity.SatHostState, Complexity.satHostEntry, Complexity.satHostRet, Complexity.satPadTM, Turing.FinTM, Turing.FinTM.controlAction, Turing.emitAction}
The finite controller hosting the six modules on one shared padded work bank:
the anchor dispatches into the recurring call sequence, and inside a module
every transition (including the halting one) is forwarded by the
output-emitting embedding, redirecting halts to the module's return state. The
clean-call contracts restore the shared scratch before each next call.
\end{definition}

\begin{lemma}[Forwarded clean seam]
\lean{Complexity.satHost_seam}
\label{Complexity.satHost_seam}
\leanok
\uses{Complexity.SatHostState, Complexity.satHostRet, Complexity.satPadCfg, Complexity.satPad_seam, Turing.Cfg, Turing.Cfg.ofWords, Turing.FinTM, Turing.emitCfg, Turing.stateWord}
\difficulty{3}
For a module with positive tape count: forwarding a padded canonical
configuration of the module (built from a state word, with a given output)
through the emitting embedding produces the host's canonical seam over the same
state word, with the supplied prefix prepended to the output; a halted module
state maps to the module's return state.
\end{lemma}

\begin{lemma}[Running one clean module in the host]
\lean{Complexity.satHost_call}
\label{Complexity.satHost_call}
\leanok
\uses{Complexity.SatHostState, Complexity.satHostEntry, Complexity.satHostRet, Complexity.satHostTM, Complexity.satHost_seam, Complexity.satPadCfg, Complexity.satPadTM, Complexity.satPad_run, Turing.Cfg.Halted, Turing.Cfg.ofWords, Turing.FinTM, Turing.MultiTapeTM.runFrom, Turing.emitCfg, Turing.emit_run, Turing.stateWord}
\difficulty{5}
Let the $i$-th module (with positive tape count), started at its own entry over
the state word of an argument, stay live before time $t$ and halt at time $t$
over the state word of a result with output $o$. Then the host, started at that
module's entry over the same argument with output prefix $p$: (i) never visits
the loop anchor before time $t$, and (ii) at time $t$ reaches the module's
return state over the state word of the result with output
$p \mathbin{+\!\!+} o$.
\end{lemma}

\begin{lemma}[Joining guarded segments]
\lean{Complexity.sat_guard_add}
\label{Complexity.sat_guard_add}
\leanok
\uses{Turing.Cfg, Turing.MultiTapeTM, Turing.MultiTapeTM.runFrom, Turing.MultiTapeTM.runFrom_add}
\difficulty{3}
Let a machine avoid a designated anchor state for $a$ steps from $c$, reach $d$
at time $a$, and avoid the anchor for $b$ further steps from $d$. Then it avoids
the anchor at every time below $a + b$ from $c$.
\end{lemma}

\begin{definition}[Clean-module contract]
\lean{Complexity.SatClean}
\label{Complexity.SatClean}
\leanok
\uses{Turing.Cfg.Halted, Turing.Cfg.ofWords, Turing.FinTM, Turing.MultiTapeTM.runFrom, Turing.stateWord}
The uniform contract for a host module $M$ with result and emission functions
and a budget $B$: $M$ has a positive tape count, and for every native input and
argument word there is a positive time within $B$ of the argument length at
which $M$, started at its entry over the state word of the argument, first halts
— in the canonical configuration over the state word of the result, having
emitted exactly the prescribed chunk, with the native input untouched.
\end{definition}

\begin{lemma}[From bridge contract to clean contract]
\lean{Complexity.satClean_stop}
\label{Complexity.satClean_stop}
\leanok
\uses{Complexity.SatClean, Complexity.satStopTM, Complexity.satStop_clean, Complexity.sat_call_first_halt, Turing.Cfg.ofWords, Turing.FinTM, Turing.MultiTapeTM.runFrom, Turing.stateWord}
\difficulty{4}
Let a machine $C$ with entry and exit states satisfy a first-positive-exit
bridge contract: on every input and argument it reaches, within budget $B$, the
exit over the state word of the result with the prescribed emission, avoiding
the exit at intermediate positive times. Then the call-stopping wrapper of $C$
satisfies the uniform clean-module contract with budget $B + 1$.
\end{lemma}

\begin{lemma}[Bridge budget bound]
\lean{Complexity.sat_bridge_bound}
\label{Complexity.sat_bridge_bound}
\leanok
\difficulty{3}
For all naturals $c, A, d, n$ and any length
$\mathrm{len} \le A\,(n+1)^d$:
\[
c\,\bigl(A\,(n+1)^d + n + \mathrm{len} + 1\bigr) + 1
\;\le\; \bigl(c\,(2A + 1) + 1\bigr)\,(n+1)^{d+1},
\]
i.e.\ the bridge overhead is bounded by one monomial of degree one higher.
\end{lemma}

\begin{lemma}[Clean install modules]
\lean{Complexity.satClean_install}
\label{Complexity.satClean_install}
\leanok
\uses{Complexity.PolyTimeComputable, Complexity.SatClean, Complexity.satClean_stop, Complexity.satStopTM, Complexity.sat_bridge_bound, Turing.FinTM.computesInTime_iff, Turing.FinTM.exists_installCallTM, Turing.MultiTapeTM.output_length_le}
\difficulty{4}
For every polynomial-time computable word function $f$ there exist a module $M$
and constants $A, d$ such that $M$ satisfies the clean-module contract with
result function $f$, empty emission, and budget $n \mapsto A\,(n+1)^d$ on the
argument length: every polynomial word computation has a clean first-halting
install module.
\end{lemma}

\begin{lemma}[Clean emit modules]
\lean{Complexity.satClean_emit}
\label{Complexity.satClean_emit}
\leanok
\uses{Complexity.PolyTimeComputable, Complexity.SatClean, Complexity.satClean_stop, Complexity.satStopTM, Complexity.sat_bridge_bound, Turing.FinTM.computesInTime_iff, Turing.FinTM.exists_emitCallTM, Turing.MultiTapeTM.output_length_le}
\difficulty{4}
For every polynomial-time computable word function $f$ there exist a module $M$
and constants $A, d$ such that $M$ satisfies the clean-module contract with the
identity as result function, emission $f$, and budget $n \mapsto A\,(n+1)^d$:
the emit-mode module preserves its argument and forwards the exact computed
chunk, with cleanup and the final halt inside the polynomial budget.
\end{lemma}

\begin{lemma}[Length of a pair encoding]
\lean{Complexity.satPair_length}
\label{Complexity.satPair_length}
\leanok
\uses{Turing.pairEncode}
\difficulty{2}
For all binary words $u$ and $v$, the pair encoding of $u$ with $v$ has length
exactly $2|u| + |v| + 2$; this charges every round request to the original
input.
\end{lemma}

\begin{lemma}[Common budget for linear requests]
\lean{Complexity.sat_request_budget}
\label{Complexity.sat_request_budget}
\leanok
\difficulty{2}
For naturals $A, d, D, n, m$ with $d \le D$ and $m + 1 \le 32\,(n+1)$:
\[
A\,(m+1)^d \;\le\; A\,\bigl(32\,(n+1)\bigr)^{D}.
\]
Every polynomial on a linearly bounded request length fits the common power in
the original input length.
\end{lemma}

\begin{lemma}[Clean module at one call site]
\lean{Complexity.satHost_clean}
\label{Complexity.satHost_clean}
\leanok
\uses{Complexity.SatClean, Complexity.satHostEntry, Complexity.satHostRet, Complexity.satHostTM, Complexity.satHost_call, Turing.Cfg.ofWords, Turing.FinTM, Turing.MultiTapeTM.runFrom, Turing.stateWord}
\difficulty{2}
Let the $i$-th module of the host satisfy the clean-module contract with result
$f$, emission $g$, and budget $B$. Then for every native input, argument, and
output prefix, there is a positive time within $B$ of the argument length during
which the host, started at that module's entry, avoids the loop anchor and then
reaches the module's return state over the state word of $f$ of the argument,
with $g$ of the argument appended to the output.
\end{lemma}

\begin{lemma}[The reduction is polynomial-time computable]
\lean{Complexity.satReduction_poly}
\label{Complexity.satReduction_poly}
\leanok
\uses{Complexity.PolyTimeComputable, Complexity.polyTimeComputable_id, Complexity.satAppendTM, Complexity.satAppend_clean, Complexity.satClean_emit, Complexity.satClean_install, Complexity.satHostEntry, Complexity.satHostRet, Complexity.satHostTM, Complexity.satHost_call, Complexity.satHost_clean, Complexity.satPair_length, Complexity.satReduction, Complexity.satReqEmit, Complexity.satReqEmit_poly, Complexity.satReqStep, Complexity.satReqStep_poly, Complexity.satReq_round, Complexity.satStreamBound_size, Complexity.satStreamEmit, Complexity.satStreamInv, Complexity.satStreamInv_start, Complexity.satStreamInv_step, Complexity.satStreamRead_word, Complexity.satStreamStart, Complexity.satStreamStart_poly, Complexity.satStreamStep, Complexity.satStreamWord, Complexity.satStream_output_identity, Complexity.sat_guard_add, Complexity.sat_pt_const, Complexity.sat_pt_pair, Complexity.sat_request_budget, Turing.Action.apply, Turing.Cfg.init, Turing.Cfg.ofWords, Turing.FinTM, Turing.FinTM.ComputesFunInTime, Turing.FinTM.ComputesInTime.mono, Turing.FinTM.bufferTape, Turing.FinTM.computesFunInTime_lengthBits, Turing.FinTM.controlAction, Turing.FinTM.exists_emitLoopTM, Turing.MultiTapeTM.initCfg, Turing.MultiTapeTM.runFrom, Turing.MultiTapeTM.runFrom_add, Turing.MultiTapeTM.runFrom_succ_eq_step, Turing.MultiTapeTM.step, Turing.pairEncode, Turing.stateWord}
\difficulty{7}
The string-level SAT-to-3SAT reduction — decode, apply the clause-splitting
transform, and serialize, with the empty-formula fallback on malformed inputs —
is polynomial-time computable: a finite host of six clean modules (startup
append and install, then the recurring pack, append, emit, and install calls)
drives the normalized streaming schedule through an emitting loop with round
budget $R(n) = n$, every request being linearly bounded in the original input.
\end{lemma}

\begin{theorem}[SAT reduces to 3SAT]
\lean{Complexity.SAT_reducible_SAT3}
\label{Complexity.SAT_reducible_SAT3}
\leanok
\uses{Complexity.SAT, Complexity.SAT3, Complexity.satReduction, Complexity.satReduction_correct, Complexity.satReduction_poly}
\difficulty{1}
$\mathrm{SAT} \le_p \mathrm{3SAT}$: there is a polynomial-time computable map on
binary strings — clause splitting with fresh variables — under which membership
in $\mathrm{SAT}$ of any string is equivalent to membership in $\mathrm{3SAT}$
of its image. [AB09, Lemma~2.14]
\end{theorem}
